\documentclass[a4paper,11pt,openany]{ltjsbook}

% Encoding / Japanese
\usepackage{cmap}
\usepackage{luatexja}
\usepackage{luatexja-fontspec}

% Math
\usepackage{amsmath,amssymb,mathtools}

% Layout
\usepackage{geometry}
\geometry{margin=25mm}
\usepackage{setspace}
\setstretch{1.12}
\usepackage{enumitem}
\setlist{nosep,leftmargin=2em}
\usepackage{booktabs,longtable,array}

% Figures
\usepackage{graphicx}
\usepackage{xcolor}
\usepackage{tikz}
\usetikzlibrary{arrows.meta,positioning,calc,fit,shapes.geometric}

% Boxes
\usepackage[most]{tcolorbox}
\tcbuselibrary{breakable,skins}

% Code
\usepackage{accsupp}
\usepackage[cache=false]{minted}
\usemintedstyle{friendly}
\newcommand{\copyablespace}{%
  \BeginAccSupp{method=escape,ActualText={\space}}%
  \phantom{\textvisiblespace}%
  \EndAccSupp{}%
}
\setminted[lean4]{
  fontsize=\small,
  breaklines=true,
  breaksymbolleft={},
  frame=single,
  framerule=0.3pt,
  framesep=0.6em,
  rulecolor=\color{gray!40},
  autogobble=false,
  obeytabs=true,
  tabsize=2,
  showspaces=true,
  space=\copyablespace
}

% Hyperlinks
\usepackage[hidelinks]{hyperref}
\usepackage[nameinlink,noabbrev]{cleveref}

% Japanese captions
\renewcommand{\figurename}{図}
\renewcommand{\tablename}{表}
\floatname{listing}{コード}

% Chapter abstract
\newenvironment{chapterabstract}{%
  \par\vspace{1.2em}{\Large\bfseries はじめに}\par\vspace{0.8em}
}{%
}

% Learning goals
\newenvironment{learninggoals}{%
  \par\vspace{1.2em}{\Large\bfseries この章の到達目標}\par\vspace{0.8em}
  \begin{itemize}
}{%
  \end{itemize}
}

% Book index entry with a reading/ASCII sort key and a reader-facing label.
% Usage: \BookIndex{reading-or-ASCII-key}{display label}
\newcommand{\BookIndex}[2]{\index{#1@#2}}

% Chapter reference that follows the target chapter's current number.
% Usage: \ChapterRef{chap:chapter-label}
\newcommand{\ChapterRef}[1]{第\ref{#1}章}

% Figure placeholder.
% Usage: \FigurePlaceholder{filename.png}{0.90}{caption}{fig:label}
\newcommand{\FigurePlaceholder}[4]{%
\begin{figure}[htbp]
\centering
\IfFileExists{figures/#1}{%
  \includegraphics[width=#2\linewidth]{figures/#1}%
}{%
  \PackageError{category-theory-lean4-for-engineers}%
    {Required figure `figures/#1' is missing}%
    {Add the referenced figure before building a release PDF.}%
}
\caption{#3}
\label{#4}
\end{figure}
}


\title{ソフトウェアエンジニアのための圏論と Lean 4 入門\\\large 仕様・同値性・安全なマイグレーションを証明する}
\author{Author Name}
\date{\today}

\begin{document}

\frontmatter
\maketitle
\tableofcontents

\mainmatter

\renewcommand{\thepart}{\arabic{part}}
\setcounter{part}{-1}
\clearpage
\refstepcounter{part}
\addcontentsline{toc}{part}{第\thepart 部\hspace{1em}地図：なぜ圏論と Lean を一緒に学ぶのか}
\begin{center}
\makebox[\textwidth][c]{\large\bfseries 第\thepart 部\quad 地図：なぜ圏論と Lean を一緒に学ぶのか}
\end{center}
\vspace{2em}
\label{part:map}

% !TEX root = ../main.tex
ソフトウェア開発では、私たちは日常的に「変更しても意味が変わっていないか」を気にしています。
関数の分割、APIレスポンスの形の変更、DBスキーマの移行、古いデータから新しい形式への変換、複数のサービスから得た情報の結合などが、その例です。
これらは一見ばらばらの作業に見えますが、「ある形のものを別の形へ変換し、その変換が期待した性質を保つか」を確認する点で共通しています。

圏論は、このような変換と合成を扱うための言語です。
対象、射、合成、恒等射、同型、関手、自然変換といった言葉は、最初は抽象的に見えます。
しかし、ソフトウェアの視点では、対象を型やデータ構造、射を関数やマイグレーション、合成をパイプライン、同型を情報を失わない相互変換、関手を構造を保った一括変換として読めます。

一方、Lean 4 は「この変換は本当に正しいのか」を機械的に確認するための道具です。
通常のテストは具体例を確認し、型は多くの不正な組み合わせを排除します。
しかし、「すべての入力についてこの性質が成り立つ」「このマイグレーションは戻せる」「このリファクタリングは振る舞いを変えない」といった主張は、テストだけでは原理的に確認しきれません。
Lean では、そのような主張を命題として書き、証明として検査できます。

この部では、本書全体の地図を示します。
圏論は研究対象としてではなく、設計、仕様、変更、同値性を整理するための道具として扱います。
Lean も本番コードを置き換えるものではなく、仕様の核心部分を小さなモデルとして表し、変更前後の意味保存を確認するための実験室として扱います。

この部を読み終えるころには、本書で扱う中心的な問いを把握できます。
その問いは、「この処理は何と何を変換しているのか」「どの性質を保つべきなのか」「それはテストで確認すべきか、型で表すべきか、Leanで証明すべきか」の三つです。

\section*{この部で扱うこと}

\begin{itemize}
\item ソフトウェアにおける正しさの種類
\item テスト、型、仕様、証明の役割分担
\item 圏論を合成と変換の言語として見る視点
\item Lean 4 を仕様確認の道具として使う視点
\item 本書で扱う実務サンプルの全体像
\end{itemize}

\section*{この部で扱わないこと}

\begin{itemize}
\item 圏論の厳密な公理化
\item Lean 4 の詳細な構文説明
\item Mathlib の高度な使い方
\item 実システム全体の完全検証
\end{itemize}

\section*{次の部への接続}

次の部では、Lean 4 の最小限の読み書きを学びます。
最初から難しい数学を形式化するのではなく、型、値、関数、命題、証明という基本要素を使って、「仕様をコードの形で書く」とはどういうことかを確認します。


% !TEX root = ../main.tex
\setcounter{chapter}{-1}
\chapter{仕様・変更・証明の地図}
\label{chap:ch00_map}
\BookIndex{ただしさ}{正しさ}
\BookIndex{しよう}{仕様}
\BookIndex{けんろん}{圏論}
\BookIndex{Lean 4}{Lean 4}

\begin{chapterabstract}
本章は、本書全体の地図です。
圏論や Lean 4 の細部に入る前に、ソフトウェア開発で何を「正しい」と呼ぶのか、変更前後で何を保存したいのか、そして正しさと保存についての主張をどのように検査可能な形へ落とすのかを整理します。
テスト、型、仕様、証明は、互いに競合する道具ではありません。
それぞれが異なる粒度で正しさを支えます。
本章では、圏論を合成と変換の言語として、Lean 4 を、仕様を検査する小さな実験室として位置づけます。
\end{chapterabstract}

\begin{learninggoals}
\item ソフトウェアにおける正しさを、実行結果、型、仕様、不変条件、変更前後の意味保存に分けて説明できます。
\item テスト、型、仕様、証明がそれぞれ何を得意とし、何を保証しないのかを区別できます。
\item 圏論の語彙を、型、関数、変換、パイプライン、マイグレーション、adapter と結びつけて読めます。
\item Lean 4 を、本番コードの代替ではなく、仕様の核を検査する小さなモデルとして位置づけられます。
\item 本書で扱うケーススタディが、どのような「保存したい性質」を扱うのかを把握できます。
\end{learninggoals}

\section{本書が扱う問い}

本書は、圏論の定義を順番に暗記する本ではありません。
また、Lean 4 の構文を網羅的に説明する本でもありません。
主題は、ソフトウェアの仕様と変更を、より正確に扱うことです。

日々の開発では、次のような判断が繰り返し現れます。

\begin{itemize}
\item 関数を分割した後も、外から見た結果は変わっていませんか。
\item API レスポンスの形を変えた後も、旧クライアントの期待を壊していませんか。
\item DB スキーマを移行した後も、必要なクエリ結果は保存されていますか。
\item 文字列や JSON を encode して decode したとき、元の情報に戻りますか。
\item 口座残高や注文合計のような不変条件は、操作後にも保たれていますか。
\end{itemize}

これらは別々の作業に見えます。
しかし、どの作業も「何かを別の形へ変換し、その変換が期待した性質を保存するか」を問うています。
この共通部分を取り出すと、圏論と Lean 4 が同じ方向を向いていることが分かります。

圏論は、変換と合成を整理するための言語です。
Lean 4 は、その言語で書いた小さな主張を機械的に確認するための道具です。
本書では、圏論を抽象数学の中だけで扱わず、Lean も学術的な証明だけに限定しません。
どちらも、変更に強い設計を考えるための実務上の補助線として使います。

図\ref{fig:ch00-book-roadmap}は、仕様と変更の問いから出発し、Lean 4、圏論、証明技法を経てケーススタディへ進む本書の順序を示します。

\FigurePlaceholder{fig-ch00-book-roadmap.png}{0.92}{本書のロードマップ}{fig:ch00-book-roadmap}

本章では、まだ本格的な Lean 証明には入りません。
最初に必要なのは、証明を書くことではなく、「何を証明すべきか」を見分けることです。
正しさの種類を混同したまま Lean を使うと、証明は書けても、実務上あまり重要でない性質だけを確認してしまいます。
一方、証明すべき核を適切に切り出せれば、短い定理でも設計レビューに大きな価値を持ちます。

\section{ソフトウェアにおける「正しさ」とは何か}
\label{sec:ch00-correctness}

\subsection{正しさは一種類ではない}

ソフトウェアの正しさは、単に「バグがない」という一語では説明しきれません。
実務では、いくつもの異なる正しさが重なっています。

\begin{table}[htbp]
\centering
\caption{正しさの分類}
\label{tab:ch00-correctness-kinds}
\begin{tabular}{p{0.20\linewidth}p{0.34\linewidth}p{0.36\linewidth}}
\toprule
種類 & 典型的な問い & 例 \\
\midrule
実行結果の正しさ & この入力で期待した出力になるか & 料金計算で、入力 1000 円に対して 1100 円が返る \\
型の正しさ & 不正な形の値を受け取らないか & ユーザー ID と金額を取り違えない \\
仕様の正しさ & どの入力でも期待する性質を満たすか & 出金後も残高が負にならない \\
変更の正しさ & 変更前後で意味が保存されるか & リファクタリング後も計算結果が同じ \\
移行の正しさ & 古い形式と新しい形式の間で情報が保存されるか & \texttt{rollback (migrate u) = u} が成り立つ \\
合成の正しさ & 部品をつないでも約束が保たれるか & adapter を通す順序を変えても業務ロジックの結果が同じ \\
\bottomrule
\end{tabular}
\end{table}

これらの正しさは、互いに置き換えられません。
ある入力に対するテストが通っても、すべての入力で不変条件が保たれるとは限りません。
型が付いていても、金額計算の式が仕様どおりであるとは限りません。
仕様書に「互換性がある」と書いてあっても、互換性の方向と情報損失の有無は別途明確にする必要があります。

「テストがあるから仕様は不要」「型があるから証明は不要」「証明したから実装全体が安全」という言い方は、いずれも保証範囲を超えています。
テスト、型、仕様、証明は、それぞれ異なる不足を補います。
どれか一つがほかを完全に置き換えるわけではありません。

\subsection{正しさは保存性として現れる}

本書では、正しさをしばしば「保存される性質」として捉えます。
この捉え方は、変更を扱うときに特に有効です。

たとえば、関数をリファクタリングするときに保存したいものは、内部構造ではなく外から見た結果です。
DB スキーマを変えるときに保存したいものは、テーブルの形そのものではなく、業務上必要なクエリの意味です。
JSON encoder と decoder で保存したいものは、文字列の見た目ではなく、データが持っていた情報です。

実務で「安全な変更」と呼ぶものの多くは、「すべてを変えない変更」ではありません。
安全な変更とは、変えてよい部分と保存すべき部分を区別し、保存すべき部分が実際に保たれる変更です。

この見方を採ると、レビューで確認すべき問いが変わります。

\begin{itemize}
\item 何が変わりましたか。
\item 何を変えてはいけませんか。
\item 変えてはいけないものは、テスト、型、仕様、証明のどれで確認するのが適切ですか。
\item その確認は、具体例に対するものですか、それともすべての場合に対するものですか。
\end{itemize}

これらの問いが、本書全体の出発点になります。

\section{テスト、型、仕様、証明の役割分担}
\label{sec:ch00-test-type-spec-proof}

\subsection{四つの道具を重ねて使う}

テスト、型、仕様、証明は、同じ「正しさ」を別の角度から支えます。

図\ref{fig:ch00-correctness-layers}は、具体例を確かめるテストから、モデル上の全入力を扱う証明まで、四つの道具が受け持つ範囲を分けて示します。

\FigurePlaceholder{fig-ch00-correctness-layers.png}{0.90}{正しさを支える四つの道具}{fig:ch00-correctness-layers}

四つの道具を混同しないため、本書ではそれぞれを次の意味で使います。
\begin{itemize}
\item \textbf{テスト}：具体的な入力例について期待結果を確認します。
\item \textbf{型}：値や関数がどの形を持つかを表し、不正な組み合わせを事前に排除します。
\item \textbf{仕様}：システムや関数が満たすべき性質を明文化します。
\item \textbf{証明}：その仕様が定義された範囲のすべての場合で成り立つことを示します。
\end{itemize}

この区別を起点に、確認したい性質ごとに適切な道具を選びます。

テストは、外部システムとの接続、実 DB の挙動、UI、ブラウザ差異、パフォーマンス、運用上の例外を確認するために重要です。
これらは、Lean の小さなモデルだけでは扱いにくい対象です。
一方、テストで選べる入力例は有限個です。
入力空間が広いとき、「選んだ例では正しい」ことと「すべての入力で正しい」ことの間には大きな差があります。

型は、テストよりも広い範囲で不正な形を排除します。
たとえば、\texttt{Nat} を使えば負の数はそもそも表現できません。
\texttt{Option User} を使えば、値が存在しない場合を型に表せます。
しかし、型が付いているからといって、関数が業務仕様を満たすとは限りません。
\texttt{Nat \to Nat} という型を持つ関数は無数にあり、その中には料金計算として正しいものも誤ったものもあります。

仕様は、期待する性質を言葉にします。
自然言語の仕様書は人間に読みやすい一方で、曖昧さも持ちます。
「適切に処理する」「互換性を保つ」「一貫性を維持する」という表現だけでは、機械は検査できません。
そのため、仕様の一部を型、関数、述語、定理へ分解する必要があります。

証明は、仕様として切り出した性質が定義された範囲で常に成り立つことを示します。
ただし、証明が保証するのはモデルの内側だけです。
Lean で \texttt{UserV1} と \texttt{UserV2} という小さなモデルの round-trip を証明しても、本番のシリアライザ、DB 制約、文字コード、外部 API の全挙動まで自動的に保証されるわけではありません。

\subsection{どの道具を使うべきか}

実務では、すべてを証明しようとするのではなく、性質に合う道具を選ぶ必要があります。
次の表は、その判断の出発点です。

\begin{table}[htbp]
\centering
\caption{性質に応じた道具の選択}
\label{tab:ch00-tool-choice}
\begin{tabular}{p{0.36\linewidth}p{0.18\linewidth}p{0.36\linewidth}}
\toprule
確認したいこと & 主な道具 & 理由 \\
\midrule
代表的なユースケースで期待結果になる & テスト & 具体的な入出力例を読みやすく残せる \\
不正な形のデータを受け取らない & 型 & そもそも表現できない値を減らせる \\
業務上守るべきルールを明文化する & 仕様 & 人間が合意すべき性質を言葉にできる \\
すべての入力で式変形が意味を変えない & 証明 & 有限個の例ではなく一般的な主張を確認できる \\
旧形式から新形式へ移して戻せる & 証明 & round-trip のような全称的性質を直接表せる \\
外部 API が実際に応答する & テスト・監視 & モデル外の環境依存を観測する必要がある \\
\bottomrule
\end{tabular}
\end{table}

Lean での証明に適しているのは、「小さく切り出せる」「仕様として重要」「具体例では網羅しにくい」「変更時に壊れると困る」という条件を満たす性質です。

\section{圏論を合成と変換の言語として見る}
\label{sec:ch00-category-language}

\subsection{圏論の最初の読み替え}

圏論という名前は抽象的に見えます。
しかし、本書では最初から高度な一般論には入りません。
まずは、ソフトウェアで慣れた言葉へ読み替えます。

\begin{table}[htbp]
\centering
\caption{圏論とソフトウェアの語彙対応}
\label{tab:ch00-category-translation}
\begin{tabular}{p{0.20\linewidth}p{0.34\linewidth}p{0.36\linewidth}}
\toprule
圏論の語彙 & ソフトウェアでの読み替え & 例 \\
\midrule
対象 & 型、データ構造、状態空間 & \texttt{UserV1}, \texttt{Order}, \texttt{Account} \\
射 & 関数、変換、処理 & \texttt{normalize}, \texttt{migrate}, \texttt{decode} \\
合成 & パイプライン、処理連結 & \texttt{validate} の後に \texttt{persist} する \\
恒等射 & 何もしない処理 & 変換層を通しても値を変えない \\
同型 & 情報を失わない相互変換 & \texttt{migrate} と \texttt{rollback} が戻し合う \\
関手 & 構造を保つ一括変換 & \texttt{List.map migrate} \\
自然変換 & 一貫した adapter & \texttt{OldResponse A} から \texttt{NewResponse A} への変換 \\
積 & 両方を持つ構造 & ペア、レコード、DTO \\
余積 & どちらか一方の構造 & \texttt{Sum}, \texttt{Either}, \texttt{Result} \\
モナド & 文脈つき処理の合成 & 失敗、状態、検証を含む処理 \\
\bottomrule
\end{tabular}
\end{table}

圏論を使うということは、すべてを数学的な表現へ言い換えることではありません。
実務上の変換を、合成できるか、戻せるか、構造を保存するか、一貫しているか、という観点で整理することです。

\subsection{変更を矢印として見る}

本書では、処理や変更をしばしば矢印として描きます。

\[
  \texttt{UserV1} \xrightarrow{\texttt{migrate}} \texttt{UserV2}
\]

この図は、単に「旧形式から新形式へ変換する」という意味だけを表すものではありません。
次の問いを引き出すための図です。

\begin{itemize}
\item 逆向きの変換はありますか。
\item 変換して戻したら元に戻りますか。
\item 逆向きにも元に戻りますか。
\item すべての値で成り立ちますか、それとも well-formed な値だけで成り立ちますか。
\item リスト全体へ \texttt{map} しても件数や順序は保たれますか。
\item adapter を通す前後で業務ロジックの結果は変わりませんか。
\end{itemize}

図\ref{fig:ch00-change-map}は、旧形式と新形式、正逆の変換、変換後にも保存したい性質を別々の要素として配置します。

\FigurePlaceholder{fig-ch00-change-map.png}{0.90}{変更と保存性}{fig:ch00-change-map}

このように捉えると、圏論の語彙は実務レビューの語彙になります。
「これは同型なのか」「これは片方向互換なのか」「この \texttt{map} は構造を保存しているのか」「この adapter は自然なのか」という問いは、抽象数学のためだけのものではありません。
安全な変更を議論するための問いでもあります。

\section{Lean 4 を仕様を検査する小さな実験室として使う}
\label{sec:ch00-lean-lab}

\subsection{Lean で何をするのか}

Lean 4 は、プログラムを書ける言語であり、命題と証明も書ける言語です。
本書では、この二面性を利用します。

たとえば、\texttt{UserV1} と \texttt{UserV2} という二つの小さな構造体を Lean で定義します。
次に、\texttt{migrate : UserV1 \to UserV2} と \texttt{rollback : UserV2 \to UserV1} を書きます。
さらに、\texttt{rollback (migrate u) = u} を定理として書きます。
Lean がその証明を受け入れれば、少なくともその小さなモデルでは、移行して戻すと元に戻ることを確認できます。

Lean で行うのは、「本番コード全体をそのまま証明する」ことではありません。
まず、仕様の核を小さな型、関数、命題として表します。
そのうえで、重要な性質がすべての入力に対して成り立つことを確認します。

次のコードは、現時点で読めなくても問題ありません。
細部の構文は、\ChapterRef{chap:ch01_reading_code}以降で一つずつ説明します。
本書を読み終えるころには、この形の主張を自分で読めるようになります。
ここでは、\texttt{migrate} と \texttt{rollback} という二つの変換があり、\texttt{rollback (migrate u) = u}（移行して戻すと元に戻る）という性質を証明したい、という点だけを確認します。

\begin{listing}[htbp]
\begin{minted}{lean4}
structure UserV1 where
  id : Nat
  name : String

deriving instance Repr for UserV1

structure UserV2 where
  id : Nat
  profileName : String

deriving instance Repr for UserV2

def migrate (u : UserV1) : UserV2 :=
  { id := u.id, profileName := u.name }

def rollback (v : UserV2) : UserV1 :=
  { id := v.id, name := v.profileName }

theorem rollback_migrate (u : UserV1) :
    rollback (migrate u) = u := by
  cases u
  rfl

theorem migrate_rollback (v : UserV2) :
    migrate (rollback v) = v := by
  cases v
  rfl
\end{minted}
\caption{\texttt{migrate}・\texttt{rollback}・round-trip}
\label{lst:ch00_roundtrip_preview}
\end{listing}

この例の要点は、Lean の構文を覚えることではありません。
\texttt{migrate} と \texttt{rollback} があるだけでは不十分であり、二つの round-trip 性質を確認して初めて「情報を失っていない」と言えます。

\subsection{証明モデルと本番実装の境界}

Lean で証明した性質は、Lean で定義したモデルの範囲で強い保証を持ちます。
現実のシステムには、モデル化していない要素が残ります。

Lean で \texttt{rollback (migrate u) = u} を証明しても、実装言語の JSON ライブラリ、DB の NULL 制約、文字コード、タイムゾーン、外部 API の仕様変更まで自動で保証されるわけではありません。
証明した範囲と、テストやレビューで確認すべき範囲を分けて書く必要があります。

したがって、本書では次の方針を採ります。

\begin{itemize}
\item 実務の仕様から、証明したい小さな性質を切り出します。
\item その性質を Lean の型、関数、述語、定理として表します。
\item 証明が通る範囲を明記します。
\item モデル外の要素は、通常のテスト、レビュー、監視と組み合わせます。
\end{itemize}

この方針によって、Lean の保証範囲を過大評価せずに済みます。
Lean は厳密な道具ですが、モデル化した範囲を超えることは証明しません。

\section{本書で扱うサンプル一覧}
\label{sec:ch00-samples}

本書では、抽象概念を先に積み上げるのではなく、ソフトウェア上の問題から出発します。
次のサンプルは、章をまたいで繰り返し登場します。

図\ref{fig:ch00-case-studies}は、各ケーススタディを、意味、情報、不変条件、クエリ結果などの保存対象へ結びつけます。

\FigurePlaceholder{fig-ch00-case-studies.png}{0.92}{ケーススタディと保存対象}{fig:ch00-case-studies}

\begin{table}[htbp]
\centering
\caption{主なケーススタディ}
\label{tab:ch00-samples}
\begin{tabular}{p{0.28\linewidth}p{0.32\linewidth}p{0.30\linewidth}}
\toprule
サンプル & 保存したい性質 & 関連する主な章 \\
\midrule
関数パイプラインのリファクタリング & 入力に対する出力が変わらない & \ChapterRef{chap:ch09_equality_refactoring}、\ChapterRef{chap:ch11_category}、\ChapterRef{chap:ch18_equations_rewriting_normalization}、\ChapterRef{chap:ch34-refactoring-as-equations} \\
\texttt{UserV1} から \texttt{UserV2} へのスキーマ移行 & round-trip によって情報が戻る & \ChapterRef{chap:ch12-isomorphism}、\ChapterRef{chap:ch27-migration-classification}、\ChapterRef{chap:ch28_user_schema_lossless_migration} \\
JSON encoder / decoder & encode して decode すると元に戻る & \ChapterRef{chap:ch12-isomorphism}、\ChapterRef{chap:ch27-migration-classification} \\
口座残高操作 & 操作後も不変条件が保たれる & \ChapterRef{chap:ch10_invariants}、\ChapterRef{chap:spec-to-lean}、\ChapterRef{chap:account-spec-proof} \\
API adapter & adapter と業務ロジックの順序を入れ替えても結果が一致する & \ChapterRef{chap:natural-transformation}、\ChapterRef{chap:function-extensionality}、\ChapterRef{chap:ch31-api-adapter-naturality} \\
DB join と整合性制約 & 移行後も必要な join 結果が保存される & \ChapterRef{chap:limits-pullbacks}、\ChapterRef{chap:db-schema-query-preservation} \\
\bottomrule
\end{tabular}
\end{table}

\subsection{関数パイプラインのリファクタリング}

リファクタリングでは、内部構造を変えても外から見た意味を変えないことが求められます。
Lean では、この主張を等式として書きます。

\[
  \texttt{newPipeline x = oldPipeline x}
\]

この形は単純に見えますが、任意の \texttt{x} についての強い主張です。
この等式が成り立てば、特定のテストケースだけでなく、モデル化されたすべての入力で意味が保存されます。

\subsection{スキーマ移行と round-trip}

\texttt{UserV1} から \texttt{UserV2} への移行では、形が変わっても情報が失われないかを確認します。
情報を完全に保つ場合は、次の二つの式を証明する必要があります。

\[
  \texttt{rollback (migrate u) = u}
\]

\[
  \texttt{migrate (rollback v) = v}
\]

片方だけが成り立つ場合は、完全同型ではなく片方向互換である可能性があります。
条件付きでしか成り立たない場合は、well-formed 条件を仕様に入れる必要があります。

\subsection{JSON encoder / decoder}

encoder と decoder も、同型や round-trip の典型例です。
ただし、実務では、空白、フィールド順、デフォルト値、省略フィールド、未知フィールド、数値表現などを考慮する必要があります。
これらの要素によって、「文字列として同じ」と「データとして同じ」が一致しないことがあります。

JSON の例では、何を同じと見なすかを仕様として切り出します。
データとして戻ればよいのか、表示文字列まで一致させる必要があるのかを区別します。

\subsection{口座残高操作}

口座操作では、入金、出金、送金を扱います。
中心となるのは不変条件です。
たとえば、「残高は負にならない」という条件を操作後にも保つ必要があります。

この性質は、実務ではテストしやすいように見えます。
しかし、境界値や失敗ケースを含むすべての入力を手作業で列挙することは困難です。
Lean では、モデル化した範囲で「すべての入力」に対する主張として扱えます。

\subsection{API adapter}

API レスポンスの wrapper を変更するとき、旧形式から新形式へ変換する adapter を作ることがあります。
このとき、単に adapter が存在するだけでは不十分です。
業務ロジックを適用してから adapter を通す場合と、adapter を通してから業務ロジックを適用する場合で結果が一致するなら、adapter は業務ロジックと干渉していないと言えます。

この性質は自然性であり、実務的には「どの型でも一貫して動く adapter の条件」と読めます。

\subsection{DB join と整合性制約}

DB スキーマ変更では、移行前後でテーブルの形が変わります。
保存したいものは、必ずしもテーブル構造そのものではありません。
業務上必要なクエリ結果が保存されること、特に join の意味が変わらないことが求められます。

このとき、単に二つのレコードをペアにするだけでは足りません。
ユーザー ID が一致すること、注文が存在すること、参照整合性が満たされることなどの制約が必要です。
Pullback は、この整合条件を満たすペアを表す構造として読めます。

\section{本書の読み方}
\label{sec:ch00-reading-guide}

\subsection{最短実務ルート}

Lean 4 と圏論の両方が初めてで、まず実務に使える感覚を得たい場合は、次の順で読みます。

\begin{enumerate}
\item \ChapterRef{chap:ch01_reading_code}から\ChapterRef{chap:ch10_invariants}で、Lean 4 における型、関数、命題、不変条件を学びます。
\item \ChapterRef{chap:ch18_equations_rewriting_normalization}から\ChapterRef{chap:function-extensionality}で、等式、場合分け、帰納法、関数外延性という証明道具を身につけます。
\item \ChapterRef{chap:ch11_category}から\ChapterRef{chap:natural-transformation}で、圏、同型、積・余積、関手、自然変換をソフトウェアの例で理解します。
\item \ChapterRef{chap:spec-to-lean}から\ChapterRef{chap:ch31-api-adapter-naturality}で、仕様証明、マイグレーション、API adapter のケーススタディを確認します。
\item 付録のチェックリストを使い、自分のプロジェクトの変更に当てはめます。
\end{enumerate}

\subsection{Lean 証明力を先に強めるルート}

Lean で仕様を書けるようになりたい場合は、次の順が適しています。

\begin{enumerate}
\item \ChapterRef{chap:ch01_reading_code}から\ChapterRef{chap:ch10_invariants}で Lean の最小構文を学びます。
\item \ChapterRef{chap:ch18_equations_rewriting_normalization}から\ChapterRef{chap:function-extensionality}で、等式、場合分け、帰納法、関数外延性を学びます。
\item \ChapterRef{chap:spec-to-lean}から\ChapterRef{chap:ch26_price_spec}で、自然言語仕様を Lean の定理へ落とす流れを確認します。
\end{enumerate}

このルートでは、圏論の抽象語彙が必要になった時点で該当箇所へ戻ってもかまいません。

\subsection{圏論を丁寧に学ぶルート}

圏論の概念そのものを丁寧に身につけたい場合も、まず\ChapterRef{chap:ch01_reading_code}から\ChapterRef{chap:ch10_invariants}、続いて\ChapterRef{chap:ch18_equations_rewriting_normalization}から\ChapterRef{chap:function-extensionality}で Lean の証明道具を確認します。
そのうえで、\ChapterRef{chap:ch11_category}から\ChapterRef{chap:monad}をゆっくり読みます。
対象、射、合成、同型、関手、自然変換、モノイド、モナドを、すべてソフトウェアの例と結びつけて理解します。
\ChapterRef{chap:relations-preorders-galois}、\ChapterRef{chap:limits-pullbacks}、\ChapterRef{chap:ch37_adjunction}以降は、これらの概念をより抽象的な構造として整理します。

どのルートでも、自作の小さな定義による概念理解と、Mathlib の \texttt{CategoryTheory} API による標準化を分けて扱います。
この分離により、Lean と圏論を同時に学ぶ負荷を減らせます。

\section{本章のまとめ}

ソフトウェアの正しさは、実行結果、型、仕様、不変条件、変更前後の意味保存など、複数の層に分かれます。

テスト、型、仕様、証明は互いに補完する道具であり、どれか一つがほかを完全に置き換えるわけではありません。

圏論は、型やデータ構造を対象、関数や変換を射、パイプラインを合成として見るための言語です。

Lean 4 は、仕様の核を小さなモデルとして表し、すべての場合に成り立つ性質を機械的に確認するための実験室です。

本書では、リファクタリング、スキーマ移行、encoder / decoder、口座操作、API adapter、DB join を通じて、「保存されるべき性質」を具体的に証明します。

% % !TEX root = ../main.tex
\chapter{型・値・関数}
\label{chap:types-functions}

\begin{chapterabstract}
本章では、Lean 4を「証明を書くためだけの特殊な道具」としてではなく、まず型と関数を扱う小さな言語として導入する。
ソフトウェアエンジニアにとって、型、値、関数は日常的な概念である。
しかしLeanでは、それらが後で仕様や証明を組み立てるための材料にもなる。
本章では、\texttt{Nat}、\texttt{Bool}、\texttt{String}、\texttt{List}、\texttt{Option}を使い、\texttt{def}による関数定義、\texttt{\#eval}による実行確認、関数合成による小さなパイプラインを扱う。
圏論の用語はまだ前面に出さないが、「型を対象、関数を変換、合成を処理の連結として見る」ための準備を行う。
\end{chapterabstract}

\begin{learninggoals}
\item Lean 4における型と値を、実務上の「値に対する約束」として読める。
\item \texttt{Nat}、\texttt{Bool}、\texttt{String}、\texttt{List}、\texttt{Option}の基本的な使い方を理解できる。
\item \texttt{def}で小さな関数を定義し、\texttt{\#eval}で動作を確認できる。
\item 関数をつないで、処理パイプラインとして読む感覚を持てる。
\item 後続章で扱う仕様、証明、圏論の語彙へ進むための最小限のLeanコードを読める。
\end{learninggoals}

\section{この章を読む理由}

ソフトウェア開発では、値の形が変わるたびに小さな約束が生まれる。
ユーザーIDは数値なのか文字列なのか。
検索結果は必ず存在するのか、それとも存在しないことがあるのか。
APIから返る値は単体なのか、リストなのか。
このような約束が曖昧なままだと、実装は動いても、仕様の読み方は不安定になる。

Leanで最初に確認するのは、まさにこの約束である。
Leanは、値に必ず型を与える。
そして関数には、入力の型と出力の型を与える。
そのため、Leanのコードを読むときは、「この処理は何を受け取り、何を返すのか」を常に意識することになる。

\begin{intuitionbox}
型は、値の置き場所というよりも、値についての約束である。
たとえば\texttt{Nat}という型は「この値は自然数として扱う」という約束であり、\texttt{Option String}という型は「文字列があるかもしれないし、ないかもしれない」という約束である。
この約束を明示することが、後で仕様を命題として書くための第一歩になる。
\end{intuitionbox}

\FigurePlaceholder{fig-ch01-types-values-functions.png}{0.90}{型・値・関数の関係}{fig:ch01-types-values-functions}

\section{ソフトウェア上の問題：値の約束が曖昧になる}

実務では、関数の名前やコメントから意図を読み取ることが多い。
しかし、名前やコメントだけでは機械的には検査されない。
たとえば、次のような状況を考える。

\begin{itemize}
\item ユーザー名が空文字列かもしれない。
\item ユーザー検索は失敗するかもしれない。
\item 正規化済みの文字列と未正規化の文字列が同じ型で混ざっている。
\item 単一の値を返す処理と、リストを返す処理が混同される。
\end{itemize}

このような曖昧さは、テストで一部は発見できる。
しかし、関数の境界で「何が入力として許されるか」「何が出力として返るか」を型に書いておくと、そもそも混ぜてはいけない値を混ぜにくくなる。
Leanでは、この型の情報が非常に中心的な役割を持つ。

\begin{softwarebox}
本章で行うことは、本番アプリケーションをLeanで書き直すことではない。
実務コードの中から、仕様の核になりそうな小さな変換を取り出し、型と関数として書いてみることである。
たとえば「入力文字列を正規化する」「IDからユーザー名を探す」「見つからない場合を\texttt{Option}で表す」といった小さな処理が対象になる。
\end{softwarebox}


\section{型注釈と\texttt{\#check}}

Leanでは、すべての式に型がある。
数値リテラル、文字列リテラル、真偽値、リスト、関数、型そのものにも型がある。
最初は、「式を見たら、その式の型を確認する」という読み方に慣れることが大切である。

型を明示したいときは、次の形で書く。

\begin{center}
\texttt{(式 : 型)}
\end{center}

この\texttt{:}は、「左の式を右の型として読む」という注釈である。
日本語で読むなら、\texttt{(3 : Nat)}は「\texttt{3}を\texttt{Nat}として扱う」と読める。
ここでの\texttt{:}は、代入ではない。
値を変換しているのでもない。
Leanに対して、その式の型を明示しているだけである。

\begin{definitionbox}
本章では、型、値、式を次のように区別する。

\begin{itemize}
\item 型は、式がどの種類の値として扱われるかを表す。
\item 値は、計算の結果として得られる具体的なものを表す。
\item 式は、値そのもの、または値を作るための書き方を表す。
\end{itemize}

たとえば\texttt{3 + 4}は式であり、評価すると\texttt{7}という値になる。
その式は\texttt{Nat}型として扱える。
\end{definitionbox}

型を確認するには、\texttt{\#check}を使う。
\texttt{\#check}は、Leanに「この名前や式の型は何か」と尋ねるコマンドである。

\begin{listing}[htbp]
\begin{minted}{lean4}
#check Nat
#check Bool
#check String

#check (3 : Nat)
#check (true : Bool)
#check ("lean" : String)

#check List Nat
#check Option String
\end{minted}
\caption{型注釈と型確認}
\label{lst:ch01-type-annotation-check}
\end{listing}

このコードでは、\texttt{Nat}、\texttt{Bool}、\texttt{String}という型名も\texttt{\#check}している。
Leanでは、型名もまた式として扱われる。
たとえば\texttt{Nat}は、自然数の値そのものではなく、自然数の型を表す名前である。

\begin{leanbox}
\texttt{\#check Nat}を実行すると、Leanは\texttt{Nat : Type}のような結果を返す。
これは「\texttt{Nat}は\texttt{Type}に属する」と読める。
\texttt{Type}は、型を集めた階層の入口である。
宇宙レベルなどの詳しい話は、本章では扱わない。
今は「型にも型がある」とだけ読めば十分である。
\end{leanbox}

\texttt{List Nat}や\texttt{Option String}は、二つの単語が並んでいるように見える。
これは、\texttt{List}や\texttt{Option}に型を渡して、新しい型を作っている。
Leanでは、型に型を渡すときにも、通常の関数適用と同じように空白で並べる。
\texttt{List Nat}は「\texttt{Nat}のリスト型」、\texttt{Option String}は「\texttt{String}があるかもしれない型」と読める。

\begin{warningbox}
\texttt{(3 : Nat)}のような型注釈は、Leanに型を教えるための書き方である。
合わない型を無理に付けることはできない。
たとえば、文字列\texttt{"lean"}を\texttt{Nat}として扱うことはできない。
型注釈は変換ではなく、確認と明示のための構文である。
\end{warningbox}

\section{基本型と値の書き方}

この節では、Leanの基本的な型と、その値の書き方を確認する。
ここで扱うのは、\texttt{Nat}、\texttt{Bool}、\texttt{String}、\texttt{List}、\texttt{Option}である。
いずれも、以後の章で繰り返し使う。

\begin{center}
\begin{tabular}{lll}
\toprule
型 & 値の例 & 読み方 \\
\midrule
\texttt{Nat} & \texttt{0}, \texttt{1}, \texttt{2 + 3} & 自然数 \\
\texttt{Bool} & \texttt{true}, \texttt{false} & 真偽値 \\
\texttt{String} & \texttt{"lean"} & 文字列 \\
\texttt{List Nat} & \texttt{[1, 2, 3]} & 自然数のリスト \\
\texttt{Option Nat} & \texttt{some 3}, \texttt{none} & 自然数がある、またはない \\
\bottomrule
\end{tabular}
\end{center}

\texttt{Nat}の値は、\texttt{0}、\texttt{1}、\texttt{2}のように書く。
足し算は\texttt{+}で書ける。
本章では自然数の基本的な計算だけを使い、整数や小数には進まない。

\texttt{Bool}の値は、\texttt{true}と\texttt{false}である。
\texttt{\&\&}は論理積を表す。
\texttt{true \&\& false}は、両方が真のときだけ真になる式である。

\texttt{String}の値は、二重引用符で囲んで書く。
文字列の連結には\texttt{++}を使える。
\texttt{"lean" ++ "4"}は、二つの文字列をつないだ文字列になる。

\texttt{List A}は、\texttt{A}型の値を並べた型である。
角括弧\texttt{[...]}を使ってリストを書く。
空リスト\texttt{[]}だけでは要素型が分からないことがあるため、必要に応じて\texttt{([] : List Nat)}のように型注釈を付ける。

\texttt{Option A}は、\texttt{A}型の値がある場合と、ない場合を表す型である。
値がある場合は\texttt{some a}、値がない場合は\texttt{none}と書く。
\texttt{none}だけでは何の\texttt{Option}か分からないことがあるため、\texttt{(none : Option Nat)}のように型注釈を付ける。

\begin{listing}[htbp]
\begin{minted}{lean4}
#eval (3 : Nat) + 4
#eval true && false
#eval "lean" ++ "4"

#eval ([1, 2, 3] : List Nat).length
#eval ([true, false] : List Bool)

#eval (some 3 : Option Nat)
#eval (none : Option Nat)
\end{minted}
\caption{基本型の値を評価する}
\label{lst:ch01-basic-values}
\end{listing}

\begin{leanbox}
\texttt{List}と\texttt{Option}は、それだけでは具体的な値の型にならない。
\texttt{List Nat}、\texttt{List String}、\texttt{Option Nat}、\texttt{Option String}のように、内側に入る型を指定して使う。
このように、型から新しい型を作るものを、本章では型を受け取る構成子として読む。
\end{leanbox}

\section{\texttt{def}による名前付き定義}

Leanでは、\texttt{def}を使って名前付きの定義を書く。
定義できるものは、値でも関数でもよい。
まずは、構文の形を正確に読むことから始める。

\begin{center}
\texttt{def 名前 (引数 : 引数の型) : 返り値の型 := 本体}
\end{center}

\texttt{:=}の左側には、定義する名前、引数、返り値の型を書く。
\texttt{:=}の右側には、その定義の本体を書く。
Leanは、本体の型が返り値の型と合っているかを検査する。

\begin{listing}[htbp]
\begin{minted}{lean4}
def answer : Nat :=
  42

def addOne (n : Nat) : Nat :=
  n + 1

def add (m : Nat) (n : Nat) : Nat :=
  m + n

def joinWithSpace (left : String) (right : String) : String :=
  left ++ " " ++ right

#check answer
#check addOne
#check add

#eval addOne 5
#eval add 2 3
#eval joinWithSpace "Lean" "4"
\end{minted}
\caption{\texttt{def}による値と関数の定義}
\label{lst:ch01-def-basics}
\end{listing}

\texttt{answer}は引数を持たない定義である。
\texttt{answer : Nat}と書かれているので、\texttt{answer}は\texttt{Nat}型の値である。

\texttt{addOne}は一つの引数を持つ関数である。
\texttt{(n : Nat)}は、引数\texttt{n}が\texttt{Nat}型であることを表す。
返り値の型も\texttt{Nat}なので、\texttt{addOne}は\texttt{Nat}を受け取り\texttt{Nat}を返す関数である。

\texttt{add}と\texttt{joinWithSpace}は、二つの引数を持つ関数である。
Leanでは、複数の引数を\texttt{(m : Nat) (n : Nat)}のように並べて書く。
関数を呼び出すときも、\texttt{add 2 3}のように、関数名の後ろに引数を空白で並べる。

\begin{warningbox}
Leanの\texttt{def}は、可変変数への代入ではない。
\texttt{def answer : Nat := 42}は、「\texttt{answer}という名前を\texttt{42}という定義に結びつける」という意味である。
あとから同じ名前に別の値を代入して更新する、という読み方はしない。
\end{warningbox}

\section{\texttt{\#eval}による評価}

\texttt{\#eval}は、Leanに式を評価させるためのコマンドである。
\texttt{def}で定義した名前は、\texttt{\#eval}の中で使える。
Leanは定義を展開し、計算できる部分を計算し、結果を表示する。

\begin{listing}[htbp]
\begin{minted}{lean4}
#eval answer
#eval addOne answer
#eval addOne (addOne 10)
#eval add (addOne 2) (addOne 3)
#eval joinWithSpace (joinWithSpace "Lean" "4") "book"
\end{minted}
\caption{定義を使った評価}
\label{lst:ch01-eval-basics}
\end{listing}

\texttt{addOne (addOne 10)}では、内側の\texttt{addOne 10}が先に使われ、その結果に外側の\texttt{addOne}が適用される。
括弧は、どの部分を先にひとまとまりとして読むかを示している。

関数適用は、Leanでは空白で書く。
\texttt{addOne 5}は、\texttt{addOne}という関数を\texttt{5}に適用する式である。
二引数の関数も同じで、\texttt{add 2 3}と書く。

\begin{leanbox}
\texttt{\#eval}は、証明ではない。
一つの式を評価して、その結果を表示するコマンドである。
後の章で扱う\texttt{theorem}や\texttt{example}は、「すべての場合に成り立つ主張」をLeanに確認させるための構文であり、\texttt{\#eval}とは役割が異なる。
\end{leanbox}

\begin{warningbox}
\texttt{\#eval}は、Leanが計算でき、結果を表示できる式に対して使う。
型を確認したいときは\texttt{\#check}、式を評価したいときは\texttt{\#eval}、という使い分けを覚えるとよい。
\end{warningbox}

\section{関数型・関数適用・無名関数}

関数にも型がある。
\texttt{A -> B}は、\texttt{A}型の値を受け取り、\texttt{B}型の値を返す関数の型である。
たとえば、\texttt{Nat -> Nat}は自然数から自然数への関数の型である。

\texttt{def addOne (n : Nat) : Nat := n + 1}と書くと、\texttt{addOne}の型は\texttt{Nat -> Nat}になる。
\texttt{\#check addOne}で確認できる。

関数を名前なしで直接書くこともできる。
そのための構文が\texttt{fun}である。

\begin{center}
\texttt{fun 引数 => 本体}
\end{center}

\texttt{fun (n : Nat) => n + 1}は、「\texttt{n}を受け取り、\texttt{n + 1}を返す関数」をその場で作る式である。
このような関数を無名関数、またはラムダ式と呼ぶ。

\begin{listing}[htbp]
\begin{minted}{lean4}
#check addOne
#check (fun (n : Nat) => n + 1)

def inc : Nat -> Nat :=
  fun n => n + 1

def applyTwice (f : Nat -> Nat) (x : Nat) : Nat :=
  f (f x)

#eval inc 4
#eval applyTwice inc 10
#eval applyTwice (fun n => n + 2) 10
\end{minted}
\caption{関数型と無名関数}
\label{lst:ch01-function-types-lambda}
\end{listing}

\texttt{inc}は、返り値の型として\texttt{Nat -> Nat}を明示している。
右辺の\texttt{fun n => n + 1}では、\texttt{n}の型を書いていない。
これは、左辺に\texttt{Nat -> Nat}と書かれているため、Leanが\texttt{n : Nat}と推論できるからである。

\texttt{applyTwice}は、関数を引数として受け取る関数である。
第一引数\texttt{f}の型は\texttt{Nat -> Nat}である。
本体\texttt{f (f x)}は、まず\texttt{x}に\texttt{f}を適用し、その結果にもう一度\texttt{f}を適用する。

\begin{leanbox}
\texttt{A -> B -> C}という型は、Leanでは右に結合して\texttt{A -> (B -> C)}と読まれる。
また、\texttt{f x y}という関数適用は、左に結合して\texttt{(f x) y}と読まれる。
この二つの規則により、複数引数の関数を自然に書ける。
\end{leanbox}

\section{関数合成と型パラメータ}

二つの関数は、出力の型と入力の型が合うときにつなげられる。
\texttt{f : A -> B}と\texttt{g : B -> C}があるなら、まず\texttt{f}を適用し、その結果に\texttt{g}を適用することで、\texttt{A -> C}型の新しい関数を作れる。
これを関数合成と呼ぶ。

Leanでは、関数合成を自分で次のように定義できる。

\begin{listing}[htbp]
\begin{minted}{lean4}
def compose {A B C : Type} (g : B -> C) (f : A -> B) : A -> C :=
  fun x => g (f x)

def surround (s : String) : String :=
  "[" ++ s ++ "]"

def mark (s : String) : String :=
  s ++ "!"

def transform : String -> String :=
  compose mark surround

#eval transform "x"
#eval compose addOne addOne 10
\end{minted}
\caption{関数合成を定義する}
\label{lst:ch01-compose-definition}
\end{listing}

\texttt{compose}の定義を順に読む。
\texttt{\{A B C : Type\}}は、\texttt{A}、\texttt{B}、\texttt{C}が型であることを表す。
波括弧\texttt{\{...\}}で囲まれているので、通常はLeanがこれらを推論する。
たとえば\texttt{compose mark surround}では、\texttt{surround : String -> String}と\texttt{mark : String -> String}から、\texttt{A}、\texttt{B}、\texttt{C}はいずれも\texttt{String}だと分かる。

\texttt{(g : B -> C)}は、\texttt{B}から\texttt{C}への関数を受け取るという意味である。
\texttt{(f : A -> B)}は、\texttt{A}から\texttt{B}への関数を受け取るという意味である。
返り値の型は\texttt{A -> C}である。
つまり、\texttt{compose}は二つの関数を受け取り、新しい関数を返す。

本体\texttt{fun x => g (f x)}では、入力\texttt{x}にまず\texttt{f}を適用し、次に\texttt{g}を適用している。
\texttt{compose g f}という順番は、「先に\texttt{f}、後で\texttt{g}」と読む。
これは数学でよく使う合成記法の順番に近い。

\begin{intuitionbox}
関数合成で最も大切なのは、型がつながることを見ることである。
\texttt{A -> B}の出力\texttt{B}が、\texttt{B -> C}の入力\texttt{B}と一致しているから、二つの関数をつなげられる。
この見方が、第5章で扱う「射の合成」につながる。
\end{intuitionbox}

\FigurePlaceholder{fig-ch01-pipeline.png}{0.90}{関数合成による文字列変換}{fig:ch01-pipeline}

\section{サンプル：小さな文字列変換パイプライン}

最後に、ここまでに出てきた構文を一つの小さな例にまとめる。
新しい考え方を増やすのではなく、型注釈、\texttt{def}、\texttt{\#eval}、\texttt{Option}、\texttt{match}、\texttt{List.map}が同じコードの中でどう並ぶかを確認する。

\begin{listing}[htbp]
\begin{minted}{lean4}
def nonempty (s : String) : Option String :=
  if s == "" then none else some s

def decorate (s : String) : String :=
  mark (surround s)

def decorateIfNonempty (s : String) : Option String :=
  match nonempty s with
  | none => none
  | some value => some (decorate value)

def decorateAll (xs : List String) : List (Option String) :=
  xs.map decorateIfNonempty

#eval decorateIfNonempty ""
#eval decorateIfNonempty "lean"
#eval decorateAll ["a", "", "b"]
\end{minted}
\caption{文字列変換パイプライン}
\label{lst:ch01-string-pipeline}
\end{listing}

\texttt{nonempty}は、空文字列の場合に\texttt{none}を返し、それ以外の場合に\texttt{some s}を返す。
返り値の型が\texttt{Option String}なので、結果は「文字列がある場合」と「ない場合」に分かれる。

\texttt{decorateIfNonempty}では、\texttt{match}を使って\texttt{Option String}を場合分けしている。
\texttt{none}の場合は\texttt{none}を返す。
\texttt{some value}の場合は、取り出した\texttt{value}に\texttt{decorate}を適用し、もう一度\texttt{some}で包む。

\texttt{decorateAll}では、\texttt{xs.map decorateIfNonempty}と書いている。
ここでの\texttt{map}は、リストの各要素に同じ関数を適用し、結果をリストとして返す操作である。
\texttt{List.map}の性質は第8章で詳しく扱う。
本章では、ドット記法\texttt{xs.map f}を「\texttt{xs}の各要素に\texttt{f}を適用する」と読めれば十分である。

\section{本章でまだ扱わないこと}

本章では、Leanを「仕様を書くための言語」として使い始めた。
ただし、まだ扱わないことも多い。

\begin{warningbox}
本章では、\texttt{Prop}、\texttt{theorem}、\texttt{example}、\texttt{rfl}、\texttt{simp}、\texttt{rw}を本格的には扱わない。
これらは次章以降で導入する。
本章の目的は、証明を急ぐことではなく、Leanの型、値、関数、実行確認を読めるようになることである。
\end{warningbox}

また、本章の文字列変換の例は、言語構文を確認するための小さな例である。
文字列処理そのものの詳細には踏み込まない。
ここでは、Leanで関数の形、場合分け、リストへの適用を確認するために、意図的に単純なモデルにしている。

\section{実務への読み替え}

本章の内容は、すぐに次のような実務上の読み替えができる。

\begin{itemize}
\item 入力検査は、\texttt{String -> Option String}のように失敗可能性を型へ出せる。
\item DTO変換は、\texttt{A -> B}という関数として小さく切り出せる。
\item 複数件の変換は、\texttt{List A -> List B}や\texttt{List A -> List (Option B)}として表せる。
\item パイプラインは、関数合成として分解して読める。
\end{itemize}

ここで大切なのは、Leanで実務システム全体を再実装することではない。
実務コードの中にある「仕様上重要な変換」を小さな関数として切り出すことである。
小さく切り出せると、次章以降で「その関数はどんな性質を満たすべきか」を命題として書けるようになる。

\section{圏論への橋渡し}

本章では、圏論という言葉をほとんど使わなかった。
しかし、すでに圏論へ進むための材料は出そろっている。

\begin{center}
\begin{tabular}{lll}
\toprule
本章で見たもの & 後の圏論での読み方 & ソフトウェアでの感覚 \\
\midrule
型 & 対象 & データの形、状態空間 \\
関数 & 射 & 変換、処理、adapter \\
関数合成 & 射の合成 & パイプライン、処理連結 \\
何もしない関数 & 恒等射 & no-op、恒等変換 \\
\texttt{List.map} & 関手の入口 & 構造を保った一括変換 \\
\texttt{Option} & 文脈つきの値 & 失敗可能性、nullable \\
\bottomrule
\end{tabular}
\end{center}

第5章では、これらを「合成できる処理の世界」として整理する。
その前に、第2章から第4章で、命題、証明、等式、不変条件を扱う。
型と関数だけでは、「この変換が正しい」とまでは言えない。
正しさを言うには、関数が満たすべき性質を命題として書く必要がある。

\section{本章のまとめ}

Leanでは、型を値に対する約束として読むとよい。

\texttt{Nat}、\texttt{Bool}、\texttt{String}、\texttt{List}、\texttt{Option}は、仕様の小さな部品を表すための入口になる。

\texttt{def}は、名前付きの値や関数を定義するための基本構文である。

\texttt{\#check}は型を確認し、\texttt{\#eval}は計算できる式を評価する。

\texttt{Option}を使うと、値が存在しない可能性を型に明示できる。

関数合成は、処理パイプラインとして読める。

本章の内容は、後の命題・証明・圏論の基礎になる。

% % !TEX root = ../main.tex
\chapter{命題・証明・仕様}
\label{chap:ch02_props_specs}

\begin{chapterabstract}
本章では、Lean 4 における命題と証明を、ソフトウェア仕様を検査可能な形にするための道具として導入する。
多くのプログラミング言語では、条件式は \texttt{true} または \texttt{false} を返す値として扱われる。
しかし仕様を扱うときには、「この関数はすべての入力でこの性質を満たす」「この変換は ID を保存する」のような一般的な主張を書きたくなる。
Lean では、そのような主張を \texttt{Prop} として表し、\texttt{example} や \texttt{theorem} によって証明する。
本章の目標は、証明の専門技術を網羅することではない。
関数に対して「何を保証したいのか」を命題として切り出し、\texttt{rfl}、\texttt{simp}、\texttt{rw} で小さく確認する感覚を得ることである。
\end{chapterabstract}

\begin{learninggoals}
\item \texttt{Bool} と \texttt{Prop} の違いを、実行時の判定と仕様上の主張の違いとして説明できる。
\item 関数が満たすべき性質を、Lean の命題として書ける。
\item \texttt{example}、\texttt{theorem}、\texttt{by} を使って小さな証明を書ける。
\item \texttt{rfl}、\texttt{simp}、\texttt{rw} の役割を、仕様確認の道具として読み分けられる。
\item 税込金額計算と ID 保存のような小さな実務仕様を、Lean の定理として表せる。
\end{learninggoals}

\section{この章を読む理由}

第1章では、Lean で型、値、関数を書く方法を見た。
たとえば、価格に税額を足す関数、入力を正規化する関数、データを別の形へ変換する関数は、\texttt{def} で定義できる。
しかし、関数が書けるだけでは仕様を扱ったことにはならない。
実務で本当に知りたいのは、しばしば次のような性質である。

\begin{itemize}
\item 税額が 0 のとき、税込金額は元の価格と同じになる。
\item DTO 変換をしても、ユーザー ID は変わらない。
\item 会員向けの送料計算では、会員の送料は必ず 0 になる。
\item 正規化関数を通してから計算しても、結果は変わらない。
\end{itemize}

これらは、1回の \texttt{\#eval} で確認できる例ではない。
「この入力では合っている」ではなく、「すべての入力でこの形の性質が成り立つ」と言いたいからである。

\begin{intuitionbox}
\texttt{Bool} は、プログラムが計算して返す真偽値である。
\texttt{Prop} は、Lean に対して「これを証明したい」と提示する主張である。
仕様を書くときには、実行時の判定だけでなく、証明対象としての主張を区別する必要がある。
\end{intuitionbox}

\FigurePlaceholder{fig-ch02-overview.png}{0.90}{第2章の概念地図：関数、命題、証明、仕様の関係}{fig:ch02-overview}

図\ref{fig:ch02-overview}は、本章の流れを表す。
まず関数がある。
次に、その関数に期待する性質を命題として書く。
最後に、その命題が正しいことを証明として与える。
この流れが、後続章で扱うリファクタリング証明、マイグレーション証明、API adapter の一貫性確認の土台になる。

\section{Bool と Prop の違い}

\subsection{Bool は計算される値である}

\texttt{Bool} は、\texttt{true} または \texttt{false} という値を持つ型である。
多くのプログラミング言語の boolean とほぼ同じように読める。
たとえば、会員なら割引対象にする、一定金額以上なら送料無料にする、といった分岐で使う。

\begin{listing}[htbp]
\begin{minted}{lean4}
def hasDiscount (member : Bool) : Bool :=
  member

#eval hasDiscount true   -- true
#eval hasDiscount false  -- false
\end{minted}
\caption{Bool は計算される真偽値である}
\label{lst:ch02_bool}
\end{listing}

コード\ref{lst:ch02_bool}では、\texttt{hasDiscount true} を評価すると \texttt{true} が得られる。
これは、実行時に計算される値である。
テストコードで特定の入力を与えて結果を確認する感覚に近い。

\begin{softwarebox}
\texttt{Bool} は、アプリケーション内の分岐や判定結果に向いている。
たとえば「このリクエストは有効か」「このユーザーは会員か」「この金額は送料無料条件を満たすか」を実行時に判定するときに使う。
\end{softwarebox}

\subsection{Prop は証明したい主張である}

一方、\texttt{Prop} は「証明対象となる命題」の型である。
値として計算することよりも、その主張が正しいことを Lean に納得させることが目的になる。

\begin{listing}[htbp]
\begin{minted}{lean4}
def DiscountSpec (member : Bool) : Prop :=
  hasDiscount member = true

theorem member_hasDiscount : DiscountSpec true := by
  rfl
\end{minted}
\caption{Prop は証明したい主張である}
\label{lst:ch02_prop}
\end{listing}

コード\ref{lst:ch02_prop}の \texttt{DiscountSpec true} は、\texttt{hasDiscount true = true} という主張である。
\texttt{theorem member\_hasDiscount} は、その主張が正しいことを Lean に示している。
ここで使っている \texttt{rfl} は、左右が定義を展開すれば同じになることを示す証明である。

\begin{definitionbox}
本書では、\textbf{命題}を「正しいかどうかを証明したい主張」として扱う。
Lean では命題の型を \texttt{Prop} と書く。
たとえば \texttt{price = addTax base tax} や \texttt{(toDto u).id = u.id} は命題である。
\end{definitionbox}

\subsection{仕様では両方を使う}

\texttt{Bool} と \texttt{Prop} は競合するものではない。
実務仕様では、しばしば両方を使う。
\texttt{Bool} は実装で判定する値として使い、その判定関数が期待どおりに振る舞うことを \texttt{Prop} として表す。

\begin{table}[htbp]
\centering
\caption{\texttt{Bool} と \texttt{Prop} の使い分け}
\label{tab:ch02-bool-prop}
\begin{tabular}{p{0.22\linewidth}p{0.34\linewidth}p{0.34\linewidth}}
\toprule
観点 & \texttt{Bool} & \texttt{Prop} \\
\midrule
役割 & 実行時に計算される真偽値 & 証明したい主張 \\
主な使い道 & 条件分岐、バリデーション、フラグ & 仕様、不変条件、事後条件 \\
確認方法 & \texttt{\#eval} やテストで値を見る & \texttt{theorem} や \texttt{example} で証明する \\
例 & \texttt{hasDiscount member} & \texttt{hasDiscount member = true} \\
\bottomrule
\end{tabular}
\end{table}

表\ref{tab:ch02-bool-prop}の違いを、最初から完全に使いこなす必要はない。
重要なのは、\texttt{Bool} は「判定結果」、\texttt{Prop} は「保証したい主張」と読むことである。

\FigurePlaceholder{fig-ch02-bool-vs-prop.png}{0.88}{\texttt{Bool} と \texttt{Prop} の違い：実行時の判定と仕様上の主張}{fig:ch02-bool-vs-prop}

\section{命題としての仕様}

\subsection{自然言語仕様を小さく切り出す}

自然言語の仕様書には、さまざまな粒度の文が混ざっている。
たとえば、次のような仕様を考える。

\begin{quote}
税込金額は、税抜価格に税額を加えた金額である。
税額が 0 の場合、税込金額は税抜価格と一致する。
\end{quote}

この文章をいきなりすべて Lean に移す必要はない。
まずは、関数と命題に分ける。

\begin{itemize}
\item 関数：税抜価格と税額から税込金額を計算する。
\item 命題：税額が 0 のとき、結果は元の価格と等しい。
\end{itemize}

\begin{listing}[htbp]
\begin{minted}{lean4}
def addTax (price tax : Nat) : Nat :=
  price + tax

def TaxSpec (price tax total : Nat) : Prop :=
  total = addTax price tax

example : TaxSpec 1000 100 1100 := by
  rfl
\end{minted}
\caption{仕様を関数と命題に分ける}
\label{lst:ch02_spec_as_prop}
\end{listing}

コード\ref{lst:ch02_spec_as_prop}では、\texttt{addTax} が計算であり、\texttt{TaxSpec} が仕様である。
\texttt{TaxSpec 1000 100 1100} は「1000 に 100 を足した税込金額は 1100 である」という命題になる。

\begin{leanbox}
Lean では、仕様をコメントとしてではなく、型のついた命題として書ける。
その命題に証明を与えると、Lean が主張と根拠の対応を検査する。
この点が、通常のコメントや設計メモとの大きな違いである。
\end{leanbox}

\subsection{仕様は関数の外側に置く}

実装関数と仕様命題を分けると、後で変更に強くなる。
関数の中に「こういうつもり」とコメントを書くだけでは、Lean はそれを検査できない。
しかし命題として外に出せば、関数を変更したときに証明が壊れる可能性がある。
壊れた証明は、仕様上の意味が変わったかもしれないというシグナルになる。

\begin{softwarebox}
実務では、すべての仕様を Lean に移す必要はない。
まずは「壊れると困る核」を選ぶ。
たとえば、ID が保存されること、件数が保存されること、手数料が 0 のとき結果が変わらないこと、round-trip が成り立つことなどである。
\end{softwarebox}

\section{example と theorem}

\subsection{example は名前のない小さな確認である}

\texttt{example} は、名前を付けずに命題と証明を書くための構文である。
本文中の小さな確認や、Lean の構文を試すときに向いている。

\begin{listing}[htbp]
\begin{minted}{lean4}
def SamePrice (before after : Nat) : Prop :=
  before = after

example : SamePrice 1200 1200 := by
  rfl
\end{minted}
\caption{example による小さな証明}
\label{lst:ch02_example}
\end{listing}

コード\ref{lst:ch02_example}は、「1200 と 1200 は同じ価格である」という命題を証明している。
この証明に名前はない。
そのため、後から別の証明で参照することはできない。

\subsection{theorem は名前付きの仕様である}

\texttt{theorem} は、名前付きの命題と証明を書くための構文である。
後から参照したい仕様、設計文書に対応させたい性質、CI で壊れたときに意味を追跡したい性質には、\texttt{theorem} を使う。

\begin{listing}[htbp]
\begin{minted}{lean4}
theorem addTax_zero (price : Nat) :
    addTax price 0 = price := by
  simp [addTax]
\end{minted}
\caption{theorem による名前付き仕様}
\label{lst:ch02_theorem}
\end{listing}

コード\ref{lst:ch02_theorem}は、「税額が 0 のとき、税込金額は元の価格と一致する」という名前付きの仕様である。
\texttt{addTax\_zero} という名前が付いているため、後の章でこの事実を使って別の証明を組み立てることができる。

\begin{definitionbox}
本書では、\texttt{example} を「その場限りの確認」、\texttt{theorem} を「名前を付けて再利用したい仕様」として使い分ける。
実務の設計文書と対応させるなら、重要な性質には \texttt{theorem} 名を付けるとよい。
\end{definitionbox}

\section{by で証明を書く}

\subsection{by は証明モードへの入口である}

Lean では、命題の後に \texttt{:= by} と書くと、証明を書くモードに入る。
ここから先では、\texttt{rfl}、\texttt{simp}、\texttt{rw}、\texttt{intro} などの証明コマンドを使って、ゴールを解消していく。

\begin{listing}[htbp]
\begin{minted}{lean4}
theorem same_price_refl (price : Nat) :
    price = price := by
  rfl
\end{minted}
\caption{by ブロックの基本形}
\label{lst:ch02_by_proof}
\end{listing}

コード\ref{lst:ch02_by_proof}では、\texttt{price = price} という命題を証明している。
これはどの \texttt{price} についても明らかに成り立つ。
Lean では、その明らかさを \texttt{rfl} という証明で伝える。

\subsection{証明はコメントではない}

証明は、コードに添える説明文ではない。
Lean が検査する構造化された根拠である。
たとえば、コメントなら間違っていてもコンパイルは通る。
しかし theorem の証明が間違っていれば、Lean は受け付けない。

\begin{warningbox}
\texttt{theorem} の名前がもっともらしくても、証明がなければ保証にはならない。
本書では原則として \texttt{sorry} を使わない。
\texttt{sorry} は「ここは未証明だが、いったん通す」という印であり、仕様保証としては未完了である。
\end{warningbox}

\section{rfl、simp、rw の最初の使い方}

本章では、証明の道具を三つだけに絞って導入する。
\texttt{rfl}、\texttt{simp}、\texttt{rw} である。
これらだけでも、小さな仕様確認には十分役に立つ。

\subsection{rfl は定義を展開すれば同じことを示す}

\texttt{rfl} は reflexivity、つまり反射律に由来する。
雑に言えば、左右が定義を展開すると同じものになるときに使える。

\begin{listing}[htbp]
\begin{minted}{lean4}
theorem addTax_def (price tax : Nat) :
    addTax price tax = price + tax := by
  rfl
\end{minted}
\caption{rfl で定義どおりの等式を証明する}
\label{lst:ch02_rfl}
\end{listing}

コード\ref{lst:ch02_rfl}では、\texttt{addTax price tax} の定義が \texttt{price + tax} なので、\texttt{rfl} で証明できる。
これは「定義どおりです」と Lean に伝えていると読める。

\subsection{simp は明らかな整理を任せる}

\texttt{simp} は、Lean が知っている単純化規則を使って式を整理する道具である。
たとえば、\texttt{x + 0 = x} のような自然数の基本的な整理や、\texttt{if true then ... else ...} のような分岐の整理に使える。

\begin{listing}[htbp]
\begin{minted}{lean4}
def chargeShipping (isMember : Bool) (shipping : Nat) : Nat :=
  if isMember then 0 else shipping

theorem member_shipping_zero (shipping : Nat) :
    chargeShipping true shipping = 0 := by
  simp [chargeShipping]
\end{minted}
\caption{simp で分岐を整理する}
\label{lst:ch02_simp}
\end{listing}

コード\ref{lst:ch02_simp}では、\texttt{isMember} が \texttt{true} の場合、\texttt{if true then 0 else shipping} は \texttt{0} に整理される。
\texttt{simp [chargeShipping]} は、\texttt{chargeShipping} の定義を展開し、明らかな分岐を整理している。

\subsection{rw は等式による書き換えである}

\texttt{rw} は rewrite の略である。
すでに分かっている等式を使って、式の一部を書き換える。
リファクタリング証明では非常によく使う。

\begin{listing}[htbp]
\begin{minted}{lean4}
def normalizePrice (price : Nat) : Nat :=
  price

theorem normalizePrice_eq (price : Nat) :
    normalizePrice price = price := by
  rfl

theorem addTax_after_normalize (price tax : Nat) :
    addTax (normalizePrice price) tax = addTax price tax := by
  rw [normalizePrice_eq]
\end{minted}
\caption{rw で既知の等式を使って書き換える}
\label{lst:ch02_rw}
\end{listing}

コード\ref{lst:ch02_rw}では、\texttt{normalizePrice\_eq} という等式を使って、\texttt{normalizePrice price} を \texttt{price} に書き換えている。
その結果、左右が同じ形になり、証明が完了する。

\begin{leanbox}
\texttt{rfl} は「定義を見れば同じ」、\texttt{simp} は「明らかな整理をまとめて行う」、\texttt{rw} は「既知の等式で置き換える」と読む。
この三つは、第3章以降の等式推論とリファクタリング証明の基本語彙になる。
\end{leanbox}

\section{サンプル：税込金額計算の単純な仕様}

\subsection{実務例を小さくする}

ここでは、税込金額計算を題材にする。
現実の税計算には、税率、端数処理、軽減税率、通貨単位、割引との順序などが関わる。
しかし本章では、最初の証明に集中するため、税率ではなく税額を直接受け取る単純なモデルにする。

\begin{softwarebox}
モデルを単純化することは、ごまかしではない。
証明したい核を小さく切り出すための設計判断である。
現実の複雑さをすべて入れる前に、まず「税額を足す」という最小仕様を Lean で確認する。
\end{softwarebox}

\subsection{関数、仕様、証明を並べる}

次のコードは、本章で学んだ要素をまとめたものである。
関数を定義し、命題として仕様を書き、名前付き定理として証明する。

\begin{listing}[htbp]
\begin{minted}{lean4}
def addTax (price tax : Nat) : Nat :=
  price + tax

def TaxSpec (price tax total : Nat) : Prop :=
  total = addTax price tax

example : TaxSpec 1000 100 1100 := by
  rfl

theorem addTax_zero (price : Nat) :
    addTax price 0 = price := by
  simp [addTax]
\end{minted}
\caption{税込金額計算の小さな仕様証明}
\label{lst:ch02_tax_spec}
\end{listing}

コード\ref{lst:ch02_tax_spec}で重要なのは、\texttt{example} と \texttt{theorem} の違いである。
\texttt{example} は、具体的な入力 \texttt{1000} と \texttt{100} についての確認である。
一方、\texttt{addTax\_zero} は任意の \texttt{price} について成り立つ主張である。
テストが具体例を確認するのに対し、定理はすべての入力に対する主張を扱う。

\subsection{ID 保存仕様を書く}

税込金額のような数値例だけでなく、データ変換の仕様も命題として書ける。
次の例では、入力用ユーザーデータを DTO に変換しても、ID が保存されることを証明する。

\begin{listing}[htbp]
\begin{minted}{lean4}
structure UserInput where
  id : Nat
  nameCode : Nat

structure UserDto where
  id : Nat
  displayCode : Nat

def toDto (u : UserInput) : UserDto :=
  { id := u.id, displayCode := u.nameCode }

def PreservesId (f : UserInput -> UserDto) : Prop :=
  forall u, (f u).id = u.id

theorem toDto_preservesId : PreservesId toDto := by
  intro u
  rfl
\end{minted}
\caption{DTO 変換で ID が保存されることを証明する}
\label{lst:ch02_id_preservation}
\end{listing}

コード\ref{lst:ch02_id_preservation}では、\texttt{PreservesId} が仕様である。
\texttt{forall u, ...} は「すべての \texttt{u} について」という意味である。
\texttt{intro u} は、任意のユーザー \texttt{u} を一つ取って証明を始める命令である。
その後、\texttt{toDto} の定義を見れば ID がそのままコピーされているので、\texttt{rfl} で証明できる。

\begin{intuitionbox}
「すべての入力で成り立つ」という仕様は、Lean では \texttt{forall} を使って表す。
証明では、任意の入力を一つ取り出し、その入力について成り立つことを示す。
この形は、後のマイグレーション証明で何度も出てくる。
\end{intuitionbox}

\FigurePlaceholder{fig-ch02-spec-proof-flow.png}{0.92}{自然言語仕様を、関数・命題・証明へ分解する流れ}{fig:ch02-spec-proof-flow}

\section{実務への読み替え}

本章で扱った Lean コードは小さい。
しかし、考え方は実務にそのまま接続できる。

\subsection{テストではなく証明に向く性質}

テストは重要である。
しかし、テストは基本的に選んだ入力例に対する確認である。
一方、Lean の定理は、命題に含まれるすべての入力を対象にする。
この違いにより、次のような性質は証明に向いている。

\begin{itemize}
\item 変換後も ID が保存される。
\item 0 を足しても結果が変わらない。
\item 会員フラグが true のとき送料は 0 になる。
\item 正規化関数が実質的に恒等関数なら、計算結果は変わらない。
\end{itemize}

ただし、すべてを証明に回す必要はない。
外部 API の応答、DB の実際の制約、パフォーマンス、UI の挙動などは、通常のテストや監視の方が適している場合も多い。

\subsection{仕様を小さく名付ける}

Lean で有用なのは、仕様に名前を付けられることである。
たとえば \texttt{toDto\_preservesId} という名前は、何を保証しているかを明確に示している。
この名前を設計文書やレビューコメントに対応させると、証明ファイルが仕様書の索引としても機能する。

\begin{softwarebox}
実務で theorem 名を付けるときは、関数名と保存したい性質を含めると読みやすい。
たとえば \texttt{migrate\_preservesId}、\texttt{rollback\_after\_migrate}、\texttt{map\_migrate\_preservesLength} のような名前である。
\end{softwarebox}

\section{よくある誤解}

\subsection{Bool を返せば仕様を書いたことになる、とは限らない}

\texttt{Bool} を返す関数は、判定を実装している。
しかし、その判定が仕様どおりであることは別の主張である。
たとえば、\texttt{isValidUser} という関数があっても、それが本当にユーザー仕様を満たしているかは別途確認する必要がある。
Lean では、その確認を \texttt{Prop} と theorem で表せる。

\subsection{具体例の証明と一般的な証明は違う}

\texttt{example : TaxSpec 1000 100 1100} は便利である。
しかし、それは特定の値についての確認である。
任意の価格についての仕様を表すには、\texttt{theorem addTax\_zero (price : Nat) : ...} のように、入力を変数として受け取る必要がある。

\subsection{simp は魔法ではない}

\texttt{simp} は強力だが、何でも証明してくれる道具ではない。
定義の展開、既知の単純化規則、明らかな分岐整理などが得意である。
一方、業務仕様そのものを Lean が自動で発見するわけではない。
何を証明したいかは、人間が命題として書く必要がある。

\begin{warningbox}
証明が短いからといって、仕様が軽いわけではない。
よい定義を置いた結果、証明が \texttt{rfl} や \texttt{simp} で終わることはよくある。
それは、仕様と実装の形がうまく揃っているという良いサインである。
\end{warningbox}

\section{本章のまとめ}

\begin{summarybox}
\begin{itemize}
\item \texttt{Bool} は実行時に計算される真偽値であり、\texttt{Prop} は証明したい命題である。
\item 仕様は、関数の外側に命題として切り出すと Lean で検査できる。
\item \texttt{example} はその場限りの確認、\texttt{theorem} は名前付きで再利用できる仕様である。
\item \texttt{by} は証明を書く入口であり、\texttt{rfl}、\texttt{simp}、\texttt{rw} は初期章で使う基本的な証明道具である。
\item 税込金額計算や DTO 変換の ID 保存のような小さな性質は、Lean の命題と証明として表せる。
\end{itemize}
\end{summarybox}

\section{次章への接続}

本章では、関数に対して仕様を命題として書き、その命題に証明を与える方法を学んだ。
次章では、この考え方をさらに進めて、リファクタリングを等式推論として扱う。
変更前後の関数が同じ結果を返すことを、\texttt{rw} や \texttt{calc} を使って段階的に示す。
これは、実務のコード変更を「意味を保存する書き換え」として扱うための最初の本格的なステップになる。

% \input{chapters/ch03_equality_refactoring}
% % !TEX root = ../main.tex
\chapter{データ構造と不変条件}
\label{chap:ch04_invariants}

\begin{chapterabstract}
本章では、Lean 4 でドメインモデルを作り、そのモデルが守るべきルールを不変条件として表す方法を学ぶ。
前章までは、関数の結果が等しいことを中心に扱った。
ここからは、値そのものが正しい形をしているか、操作の前後で重要な性質が壊れていないかを扱う。
題材として、ユーザーと口座を使う。
特に、出金処理 \texttt{withdraw} を小さくモデル化し、成功条件と失敗条件、事前条件と事後条件、そして \texttt{Subtype} による「条件を満たす値」の表し方を確認する。
\end{chapterabstract}

\begin{learninggoals}
\item \texttt{structure} を使って、Lean 4 上に小さなドメインモデルを定義できる。
\item データが守るべきルールを、\texttt{Prop} を返す述語として表現できる。
\item 事前条件、事後条件、不変条件を区別し、関数仕様として書ける。
\item \texttt{Subtype} を使って、「条件を満たす値」を型として扱える。
\item 口座残高を題材に、失敗しうる \texttt{withdraw} の仕様を Lean で表せる。
\end{learninggoals}

\section{この章を読む理由}

ソフトウェアのバグは、計算式の間違いだけから生まれるわけではない。
むしろ多くの場合、データが想定していない形になったときに起きる。
たとえば、次のような状況である。

\begin{itemize}
\item ユーザーIDが空、または未割当なのに、認可処理に流れてしまう。
\item 注文明細の合計と注文ヘッダの合計金額がずれている。
\item 口座残高が負にならない仕様のはずなのに、出金後に負の状態が作られる。
\item APIレスポンスのフィールドは存在するが、その組み合わせが業務ルールを満たしていない。
\end{itemize}

これらは、型だけで完全に防げる場合もあれば、防げない場合もある。
たとえば、Lean の \texttt{Nat} は自然数なので、負の数を値として作ることはできない。
一方で、「ユーザーIDは0ではない」「予約金額は残高以下である」「注文合計は明細合計と一致する」のような条件は、単にフィールドを並べただけでは保証されない。

\begin{intuitionbox}
データ構造は「何を持っているか」を表す。
不変条件は「その持ち物がどのような約束を満たすべきか」を表す。
\texttt{structure} が形を作り、述語が正しさを語り、定理が操作後にも約束が保たれることを確認する。
\end{intuitionbox}

\FigurePlaceholder{fig-ch04-invariant-flow.png}{0.90}{データ構造・不変条件・操作・証明の関係}{fig:ch04-invariant-flow}

\section{\texttt{structure} によるドメインモデル}

\subsection{レコードをLeanで表す}

実務のドメインモデルは、多くの場合、複数のフィールドを持つレコードとして現れる。
ユーザーならIDと名前を持つ。
口座なら所有者IDと残高を持つ。
Lean では、このようなレコード型を \texttt{structure} で定義する。

\begin{leanbox}
\texttt{structure} は、新しい型を作り、その型の値が持つフィールドを指定する構文である。
TypeScript の \texttt{interface} や、Scala/Kotlin の \texttt{data class}、Rust の \texttt{struct} に近い。
ただし Lean では、その後で命題や証明と組み合わせられる点が重要である。
\end{leanbox}

\begin{listing}[htbp]
\begin{minted}{lean4}
namespace Chapter04

structure User where
  id : Nat
  name : String
  deriving Repr, DecidableEq

structure Account where
  ownerId : Nat
  balance : Nat
  deriving Repr, DecidableEq

-- 以降の例で使うサンプル値
def sampleAccount : Account :=
  { ownerId := 1, balance := 100 }

#eval sampleAccount.balance  -- => 100
\end{minted}
\caption{ユーザーと口座の小さなドメインモデル}
\label{lst:ch04_structure_model}
\end{listing}

このコードでは、\texttt{User} と \texttt{Account} という二つの型を作っている。
\texttt{Account} の \texttt{balance} は \texttt{Nat} なので、Lean の世界では負の残高を直接作れない。
この時点で、型が一つの不正状態を排除している。

ただし、これだけで口座モデルとして十分とは限らない。
たとえば \texttt{ownerId := 0} を未割当IDとして扱う業務仕様なら、\texttt{ownerId} が \texttt{Nat} であるだけでは不十分である。
\texttt{0} も \texttt{Nat} の値だからである。

\begin{softwarebox}
レコード型は、データの「形」を表す。
しかし、多くの業務ルールは形だけでは表せない。
「数値である」ことと「有効なIDである」ことは違う。
「残高フィールドを持つ」ことと「操作後にも仕様が守られる」ことも違う。
\end{softwarebox}

\subsection{型で保証するか、述語で保証するか}

データのルールを表す方法は、大きく二つある。
第一に、型そのもので不正な値を作れなくする方法である。
第二に、値の上に述語を定義し、その値が条件を満たすことを別途証明する方法である。

\begin{center}
\begin{tabular}{lll}
\toprule
方法 & 例 & 特徴 \\
\midrule
型で表す & 残高を \texttt{Nat} にする & 負の値をそもそも作れない \\
述語で表す & \texttt{0 < ownerId} & 値とは別に条件を証明する \\
部分型で表す & \texttt{\{a // AccountValid a\}} & 値と証明を一緒に持つ \\
\bottomrule
\end{tabular}
\end{center}

本章では、この三つを順に使う。
まず \texttt{structure} で形を作る。
次に述語で不変条件を表す。
最後に \texttt{Subtype} で、条件を満たす値だけを型として扱う。

\section{述語としての不変条件}

\subsection{不変条件とは何か}

不変条件とは、データが有効である限り、常に守られていてほしい条件である。
英語では invariant と呼ばれる。
「変わらない条件」という名前だが、値がまったく変化しないという意味ではない。
値を操作しても、重要な約束が壊れないという意味である。

\begin{definitionbox}
本書では、不変条件を「データが有効な状態であるために満たすべき命題」と呼ぶ。
Lean では、多くの場合、ある型 \texttt{A} の値を受け取り \texttt{Prop} を返す関数として書く。
たとえば \texttt{AccountValid : Account -> Prop} は、「この口座が有効である」という命題を返す述語である。
\end{definitionbox}

\begin{listing}[htbp]
\begin{minted}{lean4}
def AccountValid (a : Account) : Prop :=
  0 < a.ownerId

-- 入金は残高だけを増やし、所有者IDは変えない。
def deposit (a : Account) (amount : Nat) : Account :=
  { a with balance := a.balance + amount }

-- 入金しても AccountValid は保たれる。
theorem deposit_preserves_valid
    (a : Account) (amount : Nat)
    (h : AccountValid a) :
    AccountValid (deposit a amount) := by
  simpa [deposit, AccountValid] using h
\end{minted}
\caption{口座が有効であることを述語として表す}
\label{lst:ch04_invariant_predicate}
\end{listing}

\begin{syntaxbox}
ここでは、構造体の一部だけを更新する記法と、単純化結果を既存の証明へ合わせるタクティクが出てくる。
\begin{itemize}
\item \texttt{\{ a with balance := ... \}} は、構造体 \texttt{a} をもとにして \texttt{balance} フィールドだけを変えた新しい値を作る。
\item \texttt{simpa [f, P] using h} は、定義 \texttt{f} や \texttt{P} を使ってゴールと証明 \texttt{h} を単純化し、形を合わせて閉じる。
\end{itemize}
\end{syntaxbox}

\texttt{AccountValid} は、口座の所有者IDが正であることを表している。
この例では、残高の非負性は \texttt{balance : Nat} という型で表している。
一方、所有者IDが \texttt{0} でないことは述語で表している。

\subsection{保存される条件として読む}

\texttt{deposit\_preserves\_valid} は、入金後にも \texttt{AccountValid} が保たれることを述べている。
これは実務的には、「入金処理は所有者IDの妥当性を壊さない」という仕様である。
証明は短い。
なぜなら、\texttt{deposit} は \texttt{balance} だけを更新し、\texttt{ownerId} を変更しないからである。

\begin{intuitionbox}
不変条件保存の証明は、「この処理は、壊してはいけない約束に触っていない」または「触っているが、触った後も約束を満たす」ことを確認する作業である。
\end{intuitionbox}

この見方は、後の章で何度も現れる。
スキーマ移行では「IDが保存される」。
リスト移行では「件数が保存される」。
API adapter では「業務ロジックと干渉しない」。
どの場合も、ある操作が何らかの性質を保存している。
不変条件は、その最初の練習である。

\section{事前条件・事後条件}

\subsection{操作の前に必要な条件}

不変条件と似ているが、別の概念として事前条件と事後条件がある。
事前条件は、操作を安全に実行するために、操作の前に満たしていてほしい条件である。
事後条件は、操作の後に成り立ってほしい条件である。

たとえば、出金では次のように考えられる。

\begin{itemize}
\item 事前条件：出金額は現在残高以下である。
\item 事後条件：成功後の残高に出金額を足すと、元の残高に戻る。
\item 不変条件：残高は負にならない。有効な口座IDは保たれる。
\end{itemize}

\FigurePlaceholder{fig-ch04-pre-post-invariant.png}{0.92}{事前条件・操作・事後条件・不変条件の時間的な関係}{fig:ch04-pre-post-invariant}

\begin{definitionbox}
事前条件は「関数を呼ぶ前に要求する条件」である。
事後条件は「関数を呼んだ後に保証したい条件」である。
不変条件は「操作の前後を通じて、データが有効であるために守りたい条件」である。
\end{definitionbox}

\begin{listing}[htbp]
\begin{minted}{lean4}
-- 出金可能である、という事前条件。
def CanWithdraw (a : Account) (amount : Nat) : Prop :=
  amount <= a.balance

-- 出金成功後に期待する事後条件。
def WithdrawPost
    (oldAccount newAccount : Account)
    (amount : Nat) : Prop :=
  And (newAccount.ownerId = oldAccount.ownerId)
      (newAccount.balance + amount = oldAccount.balance)
\end{minted}
\caption{出金の事前条件と事後条件}
\label{lst:ch04_pre_post}
\end{listing}

\begin{syntaxbox}
比較式は、値の大小関係を命題として書く構文である。
\begin{itemize}
\item \texttt{amount <= a.balance} は、「\texttt{amount} は \texttt{a.balance} 以下である」という命題である。
\item Lean のコードでは ASCII の \texttt{<=} と、数学記号の \texttt{≤} の両方が使われることがある。
\item 比較式は \texttt{Bool} ではなく、証明対象になる \texttt{Prop} として使われる場面が多い。
\end{itemize}
\end{syntaxbox}

\texttt{CanWithdraw a amount} は、\texttt{amount} が \texttt{a.balance} 以下であることを表す。
これは、出金を成功させるための条件である。
一方、\texttt{WithdrawPost oldAccount newAccount amount} は、出金が成功した後に期待する結果を表す。
所有者IDは変わらず、減った残高に出金額を戻すと元の残高になる。

\subsection{仕様を分解する効果}

自然言語で「出金は正しく行われる」とだけ書くと、何を証明すべきかが曖昧になる。
これを Lean で扱うには、少なくとも次のように分ける必要がある。

\begin{center}
\begin{tabular}{ll}
\toprule
自然言語の言い方 & Leanでの分解 \\
\midrule
出金できるときだけ成功する & \texttt{CanWithdraw a amount} \\
成功時は残高が減る & \texttt{WithdrawPost old new amount} \\
所有者は変わらない & \texttt{new.ownerId = old.ownerId} \\
残高は負にならない & \texttt{balance : Nat} または述語 \\
無理な出金は失敗する & \texttt{withdrawOpt a amount = none} \\
\bottomrule
\end{tabular}
\end{center}

このように分解すると、証明の対象が小さくなる。
Lean で扱いやすくなるだけでなく、仕様レビューでも議論しやすくなる。

\begin{softwarebox}
実務では、関数名やAPI名だけでは仕様にならない。
\texttt{withdraw} という名前から「何となく出金する」と読めても、成功条件、失敗条件、状態更新、保持すべき情報は別々に書く必要がある。
Leanの述語は、その分解を強制するための道具として使える。
\end{softwarebox}

\section{\texttt{Subtype} による「条件を満たす値」}

\subsection{値と証明を一緒に持つ}

ここまで、不変条件は \texttt{AccountValid a : Prop} のように、値とは別の命題として表してきた。
しかし、ときには「有効であることが証明済みの口座」だけを関数に渡したい。
そのようなときに使えるのが \texttt{Subtype} である。

Lean の表記では、\texttt{\{x : A // P x\}} は「型 \texttt{A} の値 \texttt{x} と、\texttt{P x} の証明を一緒に持つ型」を表す。
本章では、これを「条件付きの値」と読む。

\begin{definitionbox}
\texttt{Subtype} は、値とその値が条件を満たす証明を組にした型である。
\texttt{\{a : Account // AccountValid a\}} は、「有効であることが証明された \texttt{Account}」の型である。
\end{definitionbox}

\begin{listing}[htbp]
\begin{minted}{lean4}
-- AccountValid を満たす Account だけを表す型。
def ValidAccount : Type :=
  { a : Account // AccountValid a }

-- Subtype から中身の Account を取り出す。
def validAccountBalance (a : ValidAccount) : Nat :=
  a.val.balance

-- 入金後も有効な口座であることを、値と証明の組として返す。
def depositValid (a : ValidAccount) (amount : Nat) : ValidAccount :=
  Subtype.mk (deposit a.val amount)
    (deposit_preserves_valid a.val amount a.property)
\end{minted}
\caption{Subtypeで有効な口座だけを表す}
\label{lst:ch04_subtype}
\end{listing}

\texttt{a.val} は、\texttt{Subtype} の中に入っている元の値である。
\texttt{a.property} は、その値が条件を満たす証明である。
この二つを組み合わせることで、\texttt{depositValid} は「有効な口座を受け取り、有効な口座を返す」関数になっている。

\subsection{どこまで型に入れるべきか}

\texttt{Subtype} を使うと、条件を型の一部として扱える。
これは強力だが、すべての条件を \texttt{Subtype} に入れるべきとは限らない。
条件が多すぎると、関数の型が重くなり、証明も増える。
初学者の段階では、次の基準で考えるとよい。

\begin{itemize}
\item 常に守りたい条件は、型または \texttt{Subtype} に入れる候補になる。
\item 特定の操作のときだけ必要な条件は、事前条件として引数にする候補になる。
\item 実システムとの境界で初めて確認する条件は、検証関数で \texttt{Option} や \texttt{Except} を返す候補になる。
\end{itemize}

\begin{warningbox}
\texttt{Subtype} を使うと「不正な値を作れない設計」に近づく。
ただし、証明が必要になる場所も増える。
最初からすべての業務ルールを \texttt{Subtype} に押し込むのではなく、壊れると困る核となる条件から始めるとよい。
\end{warningbox}

\FigurePlaceholder{fig-ch04-subtype-gate.png}{0.88}{Subtypeを条件付きデータの入口として読む}{fig:ch04-subtype-gate}

\section{サンプル：口座残高が負にならない \texttt{withdraw}}

\subsection{失敗しうる処理としての出金}

出金は、いつでも成功する処理ではない。
残高が足りない場合は失敗する。
このような処理は、Lean では \texttt{Option Account} として表せる。
成功時は \texttt{some newAccount}、失敗時は \texttt{none} を返す。

\begin{leanbox}
\texttt{Option A} は、「値があるかもしれないし、ないかもしれない」ことを表す型である。
\texttt{some x} は成功して値 \texttt{x} がある状態、\texttt{none} は失敗または値なしの状態として読める。
\end{leanbox}

\begin{listing}[htbp]
\begin{minted}{lean4}
-- 残高が足りる場合だけ、残高を減らした口座を返す。
def withdrawOpt (a : Account) (amount : Nat) : Option Account :=
  if h : amount <= a.balance then
    some { a with balance := a.balance - amount }
  else
    none

#eval withdrawOpt sampleAccount 40  -- => some { ownerId := 1, balance := 60 }
#eval withdrawOpt sampleAccount 150  -- => none
\end{minted}
\caption{失敗しうる出金処理}
\label{lst:ch04_withdraw_option}
\end{listing}

この定義では、\texttt{amount <= a.balance} が成り立つときだけ出金に成功する。
失敗した場合は \texttt{none} を返す。
\texttt{balance} は \texttt{Nat} なので、成功時に \texttt{a.balance - amount} を計算しても、Lean の値としては自然数に収まる。
ただし、ここで重要なのは「成功する条件」を明示している点である。

\subsection{成功時の仕様を証明する}

成功時には、出金後の残高に出金額を足すと元の残高に戻ってほしい。
これは先ほど定義した \texttt{WithdrawPost} で表せる。
証明しやすくするため、まずは事前条件を証明引数として受け取る \texttt{withdrawChecked} を定義する。

\begin{listing}[htbp]
\begin{minted}{lean4}
-- 証明 h がある場合だけ呼べる出金処理。
def withdrawChecked
    (a : Account) (amount : Nat)
    (_h : CanWithdraw a amount) : Account :=
  { a with balance := a.balance - amount }

-- 成功時の事後条件。
theorem withdrawChecked_post
    (a : Account) (amount : Nat)
    (h : CanWithdraw a amount) :
    WithdrawPost a (withdrawChecked a amount h) amount := by
  unfold withdrawChecked WithdrawPost CanWithdraw at *
  constructor
  · rfl
  · exact Nat.sub_add_cancel h
\end{minted}
\caption{事前条件を受け取る出金処理と成功時仕様}
\label{lst:ch04_withdraw_checked}
\end{listing}

\begin{syntaxbox}
複数の部品から証明を組み立てる構文がここで出てくる。
\begin{itemize}
\item \texttt{constructor} は、ゴールが複数の条件から成るとき、それぞれの部分ゴールへ分解する。
\item \texttt{exact h} は、現在のゴールにぴったり合う証明 \texttt{h} を渡して終了する。
\item 行頭の \texttt{·} は、分解された複数のゴールを一つずつ書くための箇条書きのような区切りである。
\item レコードや \texttt{And} のように「部品をそろえる」場面でよく使う。
\end{itemize}
\end{syntaxbox}


この証明は二つの部分からなる。
第一に、所有者IDが変わらないことを \texttt{rfl} で示す。
第二に、\texttt{Nat.sub\_add\_cancel} を使って、\texttt{a.balance - amount + amount = a.balance} を示す。
この定理は、\texttt{amount <= a.balance} という事前条件があるときに使える。

\begin{intuitionbox}
\texttt{Nat.sub\_add\_cancel} は、出金額が残高以下なら「引いてから足すと元に戻る」と読むことができる。
ここでは数学の定理というより、出金成功時の残高計算を支える道具として使っている。
\end{intuitionbox}

\subsection{成功と失敗を分けて仕様化する}

\texttt{withdrawOpt} は \texttt{Option Account} を返すため、成功ケースと失敗ケースがある。
本章では場合分けの証明を深追いしない。
それは第13章で扱う。
ここでは、成功条件があるときに \texttt{some} が返ることと、成功条件がないときに \texttt{none} が返ることを、仕様として読む。

\begin{listing}[htbp]
\begin{minted}{lean4}
theorem withdrawOpt_success
    (a : Account) (amount : Nat)
    (h : CanWithdraw a amount) :
    withdrawOpt a amount =
      some { a with balance := a.balance - amount } := by
  unfold withdrawOpt CanWithdraw at *
  simp [h]

theorem withdrawOpt_failure
    (a : Account) (amount : Nat)
    (h : Not (CanWithdraw a amount)) :
    withdrawOpt a amount = none := by
  unfold withdrawOpt CanWithdraw at *
  simp [h]
\end{minted}
\caption{成功時・失敗時の戻り値仕様}
\label{lst:ch04_withdraw_option_spec}
\end{listing}

成功ケースと失敗ケースを分けて書くと、レビューで確認すべきことが明確になる。
「残高が足りるなら成功する」のか。
「残高が足りないなら失敗する」のか。
「失敗時に元の口座を返す」のか、それとも \texttt{none} を返すのか。
こうした設計上の選択を、型と定理で表せるようになる。

\begin{softwarebox}
実務では、失敗時の仕様が曖昧なまま実装されることがある。
例外を投げるのか、\texttt{null} を返すのか、エラーコードを返すのか、元の状態を返すのか。
Lean上で \texttt{Option} や \texttt{Except} を選ぶと、その選択が型に現れる。
\end{softwarebox}

\section{ドメインモデルと仕様を分離する}

\subsection{フィールドを増やしても仕様は別に書く}

実務のドメインモデルでは、フィールドが増える。
口座なら、通貨、凍結状態、利用上限、予約済み金額、バージョン番号などが入るかもしれない。
しかし、フィールドを増やすことと、仕様を明確にすることは別である。

たとえば、次のような拡張を考える。

\begin{itemize}
\item \texttt{frozen : Bool} を追加する。
\item \texttt{reserved : Nat} を追加する。
\item \texttt{dailyLimit : Nat} を追加する。
\end{itemize}

このとき、\texttt{structure} にフィールドを追加するだけでは、次のルールは保証されない。

\begin{itemize}
\item 凍結中の口座からは出金できない。
\item 予約済み金額は残高以下である。
\item 1日の出金額は上限以下である。
\end{itemize}

これらは、述語や事前条件として別に書く必要がある。
この分離は重要である。
構造体はデータの容器であり、述語はその容器に入っている値の意味を語る。

\begin{warningbox}
フィールドがあることは、仕様があることではない。
\texttt{balance : Nat} は負の値を排除するが、「十分な残高があるときだけ出金できる」ことまでは自動で保証しない。
操作の仕様は、関数定義と定理として別に書く必要がある。
\end{warningbox}

\subsection{型・述語・定理の役割分担}

本章で使った道具を整理すると、次のようになる。

\begin{center}
\begin{tabular}{lll}
\toprule
道具 & Leanでの形 & 役割 \\
\midrule
データ構造 & \texttt{structure Account} & フィールドを持つ型を作る \\
不変条件 & \texttt{AccountValid : Account -> Prop} & 有効な状態を表す \\
事前条件 & \texttt{CanWithdraw a amount} & 操作前に必要な条件を表す \\
事後条件 & \texttt{WithdrawPost old new amount} & 操作後に保証したい条件を表す \\
条件付き値 & \texttt{\{a // AccountValid a\}} & 値と証明を一緒に持つ \\
仕様証明 & \texttt{theorem ...} & 操作が条件を満たすことを示す \\
\bottomrule
\end{tabular}
\end{center}

この役割分担を意識すると、Leanコードが「実装の写し」ではなく「仕様の模型」になる。
実システムのコードをすべてLeanへ移す必要はない。
壊れると困る性質を小さく切り出し、型・関数・述語・定理へ分解すればよい。

\section{実務への読み替え}

本章の考え方は、口座処理以外にも使える。
たとえば、次のように読み替えられる。

\begin{center}
\begin{tabular}{lll}
\toprule
実務例 & 不変条件の例 & 操作の例 \\
\midrule
ユーザー & IDが有効である & プロフィール更新 \\
注文 & 合計額が明細合計と一致する & 明細追加・削除 \\
在庫 & 予約数が在庫数以下である & 引当・キャンセル \\
権限 & ロールが許可集合に含まれる & 権限追加・削除 \\
APIレスポンス & 成功時だけ本文を持つ & wrapper変換 \\
\bottomrule
\end{tabular}
\end{center}

重要なのは、最初から完全な業務モデルを作らないことである。
たとえば注文モデルをいきなり完全に形式化しようとすると、税、割引、配送、返品、在庫、キャンペーン、決済状態が絡む。
初めて Lean で仕様を書くときは、その中から一つだけ選ぶ。
「合計額が保存される」「IDが変わらない」「失敗時に状態が変わらない」のように、小さな性質に切り出す。

\begin{softwarebox}
Leanで証明する価値が高いのは、テストケースが増えやすいが、実は一般的な性質として言える部分である。
不変条件、事前条件、事後条件は、その候補を見つけるためのチェックリストになる。
\end{softwarebox}

\section{よくある誤解}

\subsection{「型があれば不変条件はいらない」}

型は強力だが、すべての業務ルールを自動的に保証するわけではない。
\texttt{Nat} は負の残高を排除できる。
しかし、「出金額は残高以下でなければならない」という操作上の条件は、別に書く必要がある。

\subsection{「不変条件はコメントに書けばよい」}

コメントは人間にとって有用だが、Leanはコメントを検査しない。
\texttt{AccountValid : Account -> Prop} のように命題として書くと、その条件を定理の前提や結論として使える。
仕様がレビュー可能なだけでなく、機械的に検査可能な形になる。

\subsection{「\texttt{Subtype} を使えばすべて安全になる」}

\texttt{Subtype} は値と証明を一緒に持つため、不正な値を渡しにくくできる。
しかし、証明を作る責任が消えるわけではない。
また、実システムから入力されるJSONやDB行は、最初から \texttt{Subtype} として手に入るわけではない。
境界では検証関数が必要になる。

\begin{warningbox}
\texttt{Subtype} は「入力のバリデーションを不要にする魔法」ではない。
外部入力を受け取る場所では、普通の値を検査し、条件を満たす場合だけ \texttt{Subtype} に変換する、という段階が必要である。
\end{warningbox}

\section{本章のまとめ}

\texttt{structure} は、ドメインモデルの形をLean上に作るための基本道具である。

不変条件は、値が有効な状態であるために守るべき条件であり、Leanでは \texttt{A -> Prop} の形で表せる。

事前条件は操作の前に必要な条件、事後条件は操作の後に保証したい条件である。

\texttt{Subtype} を使うと、値と「条件を満たす証明」を一緒に持てる。

\texttt{withdraw} のような失敗しうる処理は、\texttt{Option} を使って成功ケースと失敗ケースを明示できる。

仕様証明では、実装全体ではなく、壊れると困る小さな性質を切り出すことが重要である。

本章で、型・関数・命題に加えて、述語、不変条件、事前条件、事後条件がそろった。
次の部では、これらを圏論の言葉で見直す。
型は対象、関数は射、操作の合成はパイプラインとして読める。
本章の \texttt{Account} や \texttt{withdraw} は、後の仕様証明やマイグレーション証明で再び登場する考え方の入口である。

% \input{chapters/ch05_category}
% \input{chapters/ch06_isomorphism}
% ...

\backmatter
\bibliographystyle{plain}
\bibliography{references}

\end{document}
